% 本文件由 LaTeX 排版 Agent 根据有效 translation.json 自动生成。
% author=John Chrysostom; work_id=1922; title=论司铎职

\chapter{论司铎职}

% structure: p0001 source=h1 target=omitted reason=page-title-already-rendered-by-parent

\Emph{卷一}\newline{}\Emph{卷二}\newline{}\Emph{卷三}\newline{}\Emph{卷四}\newline{}\Emph{卷五}\newline{}\Emph{卷六}\par

\SourceRecord{Translated by W.R.W. Stephens.；Nicene and Post-Nicene Fathers, First Series；Vol. 9.；Edited by Philip Schaff.；Buffalo, NY: Christian Literature Publishing Co.,；1889.；<http://www.newadvent.org/fathers/1922.htm>.}

% structure: p0001 source=h1 target=section reason=page-title

\section{论司铎职（卷一）}

1. 我曾有许多真诚而忠实的朋友，他们明了友谊的法则，并忠实地遵守；但在这一大群人中，有一位在对我情谊上远超其余，他竭力超越他们，正如他们超越一般相识者一样。他是常在我身边的一位；因为我们一同求学，拜同样的老师。我们对所钻研的学问怀着同样的热忱和勤奋，同一种境遇所激起的炽热渴望在我们两人中同样强烈。因为不仅在学校时，而且离开学校后，当我们必须考虑选择何种生活方式最为适宜时，我们发现彼此心意相同。\par

2. 除此之外，还有其他事物也维系并保持这份和谐，使其牢不可破。因为就我们祖国的伟大而言，无人能自夸胜过对方；我也未尝因富有而骄傲，他亦未尝因贫乏而困窘；我们的家境正如我们的志趣一样相近。我们的家庭也门第相当，一切皆与我们的性情相合。\par

3. 但当我们应当追求隐修士的幸福生活与真正的哲学时，我们的天平不再平衡，他的秤盘高高升起，而我仍缠绕于此世的欲望中，将我的秤盘压低，用青年所易沉溺的幻想将其加重。此后，我们的友谊固然如前坚固，但交往却中断了；因为兴趣不同的人不可能长久相处。然而，当我开始稍稍从世俗的洪流中浮现时，他便张开双臂迎接我；但即便如此，我们仍不能保持早先的平等：因他在时间上领先于我，又表现出极大的热忱，他再度超越了我的水平，飞升至高处。\par

4. 然而，他是一位善人，珍视我的友谊，便从其余（弟兄）中分别出来，与我共度全部时光。这本是他早就渴望的，但如我所言，因我的轻浮而受阻。因为一个出入法庭、为舞台的兴奋而心绪不宁的人，不可能常与一个埋头书本、足迹不入市场的人为伴。因此，当障碍消除，他将我带入与他相同的生活状态后，便尽情释放了他长久以来压抑的渴望。他片刻也不愿离开我，执意要求我们各自离开自己的家，共同居住——事实上，他说服了我，此事便着手进行。\par

5. 但我母亲的不断哀哭阻止了我成全他的美意，或者不如说阻止了我从他手中领受这份恩惠。因为她察觉我有此意向后，便带我进入她的私室，坐在她生我时所躺的床边，泪如泉涌，并说出比哭泣更令人怜悯的话语，其哀诉如下：\Quote{我的孩子，上天并未恩准我长久享有你父亲德行的益处。因为他的死亡紧随我生你时的阵痛而来，使你成为孤儿，使我未及适时即成为寡妇，面对寡居的全部恐怖——只有亲历者才能切实理解。因为任何言语都无法充分描述一个年轻女子的困境：她刚离开父家，未经世事，却突然被巨大的悲痛撕裂，被迫承担起超越其年龄与性别的重担。她必须纠正仆人的懒散，提防他们的诡诈，抵御亲属的图谋，勇敢承受收税人的威胁以及税率征收的苛刻。若逝者留下一个孩子，即使是女儿，也会给母亲带来极大的忧虑，虽则免去了许多开销与恐惧；但若是儿子，则使她每日充满万千焦虑与无数忧惧，更不必说若要将他培养成自由人所需付出的巨大开销。然而，这些事无一能促使我再次结婚，或引入第二个丈夫进入你父亲的房屋：我坚守原状，身处风暴与喧嚣之中，未曾躲避寡居的铁炉。我首要的帮助确实来自上天的恩宠；但在那些可怕的试炼中，不断的注视你的面容，并在你身上保存逝者活生生的形象——一个相当神似的形象——对我而言并非小小的安慰。}\par

\Quote{因此，即使你在婴孩时期，还不会说话——那是孩子给父母最大欢愉的时期——你也给了我极大的安慰。诚然，你不能抱怨说，虽然我勇敢地承受了我的寡居，但减少了你的遗产，我知道这是许多不幸成为孤儿的人的命运。因为，除了保全全部遗产完整无损外，凡是为了让你获得体面地位所需的费用，我都毫不吝惜，为此动用了我自己的一些财产和我结婚的嫁妆。然而，不要以为我说这些是在责备你；只是作为这一切恩惠的回报，我恳求你一件事：不要使我陷入第二次寡居；也不要唤醒现已安息了的悲伤：等待我的死亡吧，也许不久我将离去。年轻人确实期盼着遥远的老年；但我们这些已老迈的人，除了死亡，无可期待。那么，当你把我的尸体安葬入土，将我的骨殖与你父亲的合在一起后，就去远航吧，驶向你愿去的任何海域：那时将没有人阻拦你；但只要我还活着，请满足于我同住。我恳求你，不要无谓地对抗天主，无缘无故地将我这从未亏待过你的人卷入如此巨大的不幸中。因为，如果你有任何理由抱怨我使你卷入世俗的忧虑，强迫你处理事务，那么不要被对自然法则、教养或习俗的任何敬意所束缚，就视我为仇敌而逃离我吧；但相反，如果我所做的一切都是为了为你此生的旅程提供闲暇，那么至少让这一纽带——即使没有别的——使你留在我身边。因为，即使你说有一万人爱你，也没有人能给你如此多的自由享受，因为没有人像你这样关心你的福祉。}\par

6. 这些话，以及更多，是我母亲对我说的，我转告给了那位高贵的青年。但他非但没有因这些话而气馁，反而更加迫切地提出同样的请求。当时我们正处于这种情形：他不断恳求，而我则拒绝同意。突然有消息传来，说我们即将被提升到主教的高位，这使我们两人都感到不安。我听到这个传言后，立刻充满了恐惧和困惑：恐惧的是我可能会被迫接受这一职位，困惑的是我常自问，这些人心中怎么会产生关于我们的这种想法；因为我审视自己，发现我根本不配得到这样的荣誉。但是，那位高贵的青年私下找到我，和我商议这些事，就像对一个完全不知情的人一样，恳求我们这次也像从前一样，采取同样的行动和决策：因为他愿意跟随我，无论我走哪条路，是设法逃避还是服从被选。察觉他的热忱，并考虑到，如果因我自己的软弱而使基督的羊群失去一位如此善良且完全有能力管理众多灵魂的青年，那将会使整个教会遭受损失，我便没有向他透露我的打算，尽管我此前从未向他隐瞒过任何计划。我告诉他，最好将此事推迟到另一个时机再做决定，因为眼下并不紧迫，并以此说服他放下此事，同时鼓励他，如果这样的事真的发生在我们身上，我会和他同心。但不久之后，当那位要祝圣我们的人到来时，我隐藏了自己，而\Person{巴西略}对此毫不知情，被以另一个借口带走，被迫接受了这轭，他因我曾向他许下的诺言，指望我一定会跟随，或者更确切地说，他以为是他跟随着我。因为当时在场的某些人，看见他因被捉拿而恼怒，便欺骗他，喊道：\Quote{多么奇怪，一个通常被认为脾气更暴躁的人（指我）已经非常温和地顺从了诸位教父的判决，而他却是一个被认为更明智、更温和的人，却表现出暴躁和自负，不服管教，倔强且好争辩。}他屈服于这些劝告，后来得知我逃脱了被选，就心情沉重地来找我，坐在我旁边想要说话，但因内心痛苦、无法用言语表达他所遭受的强迫而受阻。他刚一张口，悲伤就截断了他的话，使他无法说出来。我看到他泪流满面、激动不安，又知道原因，便高兴地笑了，握住他的右手，强吻了他，并赞颂天主，因为我的计划如此成功地实现了，正如我始终祈祷的那样。但当他看到我欣喜若狂、满面春风，并明白自己被我欺骗了时，他更加烦恼和痛苦。\par

7. 当他稍微从这种心神激动中恢复过来时，他说：\Quote{如果你拒绝了分配给你的部分，并且不再关心我——我确实不知道为什么——你至少应该顾及自己的名声；但事实上，你已使众人开口，世人都说你因贪图虚妄的光荣而辞去了这职务，而且没有人能使你摆脱这项指责。至于我，我无法忍受上市场；每天有太多人上前来责骂我。因为在城中任何地方遇见我时，我所有的挚友都把我拉到一旁，把大部分责任推到我身上。知道你的意图，他们说：\InnerQuote{你的事没有一件能瞒过他，你不该隐瞒，而该告诉我们，我们本不难想出一个计划来捕获他。}但我太羞愧、太惶恐，不敢告诉他们我长久以来不知你在策划这个诡计，免得他们说我们的友谊只是虚假的。因为即使情况如此——事实上也确实如此，你也不会否认你对我所做的一切——但最好还是向外界和那些对我们有中等看法的人隐藏我们的不幸。我羞于告诉他们真相以及我们之间的真实情况，将来我不得不保持沉默，低头看地，避开我所遇见的人。因为如果我逃脱了前一项指控，我又不得不因说谎而受审判。因为他们决不会相信我，当我说你把\Person{巴西略}列入那些不被允许知道你的秘密事务的人当中。然而，对此我并不太在意，既然这已合你心意，但我们如何忍受未来的耻辱呢？因为有人指责你骄傲，有人指责你虚荣；那些对我们较为宽容的指责者，将这两项过失都归咎于我们，还说我们侮辱了那些尊敬我们的人，即使他们遭受了更大的侮辱也是活该，因为他们放过了那么多杰出人物，却提拔不过是昨日还沉浸于世俗利益的年轻人，达到他们从未梦想过的尊位，以便他们能暂时皱眉、穿上深色衣服、摆出一副严肃的面容。那些从青年初期到老年一直勤修克苦的人，如今却要听命于连治理这职务的法律都未曾听过的年轻人。我不断被说出这类话乃至更糟话的人所困扰，不知如何回答他们；但我恳求你告诉我：我并不认为你逃走并招致这些杰出人物如此仇恨是没有理由或轻率的，而是你的决定是经过深思熟虑的；由此我也猜想你已准备好辩护之词。那么请告诉我，我是否能对那些指责我们的人提出任何合理辩解。}\par

因为我并不要求你为我所遭受的错待给出任何交代，也不要求你为所行的欺骗或背信弃义，或你过去从我身上得到的好处交代。因为我可以说将我的生命交在你手中，而你却待我如此诡诈，仿佛你的职责就是防备一个敌人。然而，如果你知道我们这项决定是有益的，你就不应逃避这益处；如果相反是有害的，你就也应该使我免于损失，正如你常说你尊崇我胜过他人。但你所做的一切都使我落入网罗：你本不需要向一个向来对你言语行动极为真诚坦率的人施用诡诈和虚伪。尽管如此，如我所说，我现在并不指责你这些事，也不因你中断了那些我们常常从中获得不少乐趣和益处的会谈，而使我陷入孤独境地来责备你；但这一切我都放过，默默温柔地承受，不是因为你得罪我时行事温柔，而是因为自从我珍视你的友谊那日起，我就为自己定下规矩：无论你使我遭受何种痛苦，我绝不强迫你非道歉不可。因为你自己知道，你确实给我造成了不小的损失，至少如果你还记得我们常说的话，以及外界关于我们所说的话：我们同心合意并彼此友谊守护是一大收获。的确，每个人都说我们的和睦会给我们以外许多人带来不小的益处；然而就我而言，我从未看出这如何能有益于他人；但我确实说过，我们至少能从中获得这个益处：那些想与我们争战的人会发现我们难以制服。我不断提醒你这些事，说道：\Quote{这时代是残酷的，心怀诡计的人很多，真诚的爱已不复存在，嫉妒的致命毒害取而代之；我们行走在网罗之中，在城垛边缘；那些乐于见到我们遭遇不幸的人，若我们遭遇任何不幸，他们众多，从四面八方围困我们；然而却没有人与我们同哀，或至少这等人数目容易数算。要小心，别因分离而招致许多嘲笑，以及比嘲笑更严重的损害。兄弟互助如同坚固的城，如同设防的国。不要解散这份真诚的亲密，也不要拆毁这堡垒。}我不断地说这样的话和更多的话，并不是因为我曾怀疑过这类事，而是认为你对我完全忠诚，我做了多余的预防，希望保护一个本来健康的人；但不知不觉地，我似乎是在给一个病人服药；即便如此，我也未能有幸带来任何好处，我的过度预虑也一无所获。因为你完全抛弃了所有这些考虑，毫不顾念，将我像没有压舱物的船一样抛在未知的海洋上飘荡，毫不理会我必须面对的那些汹涌波涛。因为如果我命定要遭受诽谤、嘲笑或任何其他侮辱或威胁（这类事情必然频繁发生），我该逃向谁寻求庇护？向谁诉说我的痛苦？谁会愿意救助我，击退我的攻击者，阻止他们的攻击？谁会安慰我，使我准备好承受将来可能面对的粗俗戏谑？没有人，因为你远离这场可怕的争斗，甚至听不到我的呼喊。那么你看到你造成了多大的祸害吗？现在你下了手，你意识到你造成了怎样致命的创伤吗？但让这一切过去吧：因为过去无法挽回，也无法在无路可走的困难中找到道路。我该对外界说些什么？我该如何为他们的指控辩护？\par

\Person{金口若望}：\Quote{放心吧，我回答说，因为我不仅准备在这些事上为自己辩护，而且也将尽我所能为你没有要我交代的那些事做出解释。的确，如果你愿意，我将把它们作为我辩护的起点。因为如果我只考虑外界的称赞，尽力使他们停止指控，却未能说服我最亲爱的朋友我没有冤枉他，反而以大于他为我所显示的热心的冷漠对待他——他对我如此宽容，甚至不指控他认为从我这里所受的错待，并且为了我的缘故不顾自己的利益——那么这将是我愚蠢得奇怪的事情。}\par

我究竟对你们做了什麼错事，竟使我决定从这一点上开始踏入辩护的海洋？是我欺骗了你们，隐瞒了我的意图吗？然而我这样做，是为了那被欺骗的你和那些我藉此欺骗而交给你的人的益处。因为如果欺骗本身就是恶，并且绝不可使用它，那么我愿意承受你们所要求的任何惩罚；或者更确切地说，既然你们永远不会忍心惩罚我，我将自己承受与法官在控告者定罪恶人时所宣布的同样判决。但如果欺骗并非总是有害，而是根据使用者的意图变成好或坏，那么你们就必须停止抱怨被欺骗，并证明它是出于恶意而针对你们设计的；只要这个证明缺失，那些愿意谨慎行事的人，不仅应当避免责备和控告，甚至还应当友好地接待那欺骗者。因为一个适时且动机正直的欺骗有如此大的好处，以致许多人常常因不行欺诈而受到惩罚。如果你们考察从远古时代起享有最高声誉的将领的历史，就会发现他们的大部分胜利都是通过计谋取得的，并且这些胜利比在公开战斗中取得的更受称赞。因为后者在战争中消耗更多的金钱和人力，以致他们从胜利中一无所获，反而像战败者一样在军队牺牲和资金耗尽方面遭受同样多的痛苦。此外，他们甚至不能完全享受胜利的荣耀；因为那些阵亡者获得了不小的一部分，因为他们精神上是胜利的，失败只是身体上的：因此，如果他们受伤时不会倒下，死亡没有来结束他们的辛劳，那么他们的英勇将永无止境。但那个能通过计谋取得胜利的人，不仅使敌人遭受灾难，而且使之陷入可笑境地。再者，在另一种情况下，双方同样获得授予勇气的荣誉，而在这里，他们却不能同样获得授予智慧的荣誉，奖品完全归于胜利者；另一个同样重要的点是，他们为国家保留了胜利的纯粹喜悦；因为丰富的资源和众多的人不同于心智的力量：前者如果在战争中持续使用，必然会耗尽，并辜负拥有它们的人，而智慧的本质却是越运用越增长。不仅是在战争中，而且在和平中，欺骗的需要也可发现，不仅涉及国家事务，而且涉及私人生活，如丈夫与妻子、妻子与丈夫、儿子与父亲、朋友与朋友之间，以及子女与父母之间。因为撒乌耳的女儿若不是欺骗了她的父亲，就无法将她的丈夫从撒乌耳手中救出。她的兄弟为了拯救那个她曾救出、却又再次陷入危险的人，使用了和妻子相同的武器。\BibleRef{撒上 20:11}\par

\Person{巴西略}：但这些情况都不适用于我，因为我并非敌人，也不是那些试图伤害你们的人之一，恰恰相反。因为我将自己的所有利益都托付给你们的判断，并且每当你吩咐时，我总是遵从。\par

\Person{金口圣若望}：但是，我卓越而优秀的朋友，这正是我事先说明使用这种欺骗是一件好事的理由，不仅在战争中、对待敌人时如此，而且在和平时期、对待我们最亲爱的朋友时也是如此。为了证明它不仅对欺骗者有益，而且对被欺骗者也有益：如果你去找任何一位医生，问他们如何使病人摆脱疾病，他们会告诉你，他们不仅依靠自己的专业技能，有时还借助欺骗来引导病人恢复健康，并将从中获得的帮助与他们的医术相结合。因为当病人的任性妄为和病痛的顽固性挫败了医生的计谋时，就有必要戴上欺骗的面具，以便像在舞台上一样，他们能够掩盖实际发生的情况。但是，如果你愿意，我将给你讲述一个我从医者之子那里听来的许多计谋中的一个例子。曾经有一个人突然被一场极严重的热病侵袭；灼热加剧，病人拒绝能够退热的药物，却渴望喝一口纯酒，热烈地恳求所有靠近他的人给他酒，让他满足这个致命的渴望——我说致命，因为如果有人满足了这个请求，他不仅会使热病加剧，而且会使这个不幸的人发狂。于是，专业技能受挫，资源耗尽，完全无用时，计谋登场，并以我将要讲述的方式展现了它的力量。因为医生取了一个刚从炉中取出的陶杯，将其浸在酒中，然后空着拿出来，装满水，并吩咐用窗帘遮暗病人躺卧的房间，使光线无法揭穿这个诡计，然后假装杯子里装的是纯酒，递给他喝。那人还未将杯子拿到手中，就被气味欺骗了，没有等检查给他的是什么，就被气味说服，并被黑暗欺骗，急切地一口气喝干了饮料，喝饱之后，立刻摆脱了窒息感，逃离了迫在眉睫的危险。你看到欺骗的好处了吗？如果有人要把医生所有的诡计都算出来，这清单会没完没了。而且，不仅医治身体的人，而且照顾灵魂疾病的人也经常使用这种疗法。因此，蒙福的\Person{保禄}吸引了那些犹太人群众：\BibleRef{宗 21:26} 为了这个目的，他给\Person{弟茂德}行了割损，尽管他在信中警告迦拉达人 \BibleRef{迦 5:2}，基督不会有益于那些受割损的人。因此，他服从了法律，尽管他在接受了在基督内的信德之后，将来自法律的正义视为损失。\BibleRef{斐 3:7} 因为欺骗的价值很大，只要它不是以恶意的意图引入的。事实上，这类行动不应被称为欺骗，而应被称为一种良好的管理、智慧和技能，能够在资源不足时找到方法，并弥补心智的缺陷。因为我不会称\Person{丕乃哈斯}为杀人犯，尽管他用一击杀了两个人：\BibleRef{户 25:7} 也不会称\Person{厄里亚}在杀了那一百名士兵和军官之后为杀人犯，\BibleRef{列下 1:9} 以及他因消灭那些向魔鬼献祭的人而造成的血河。\BibleRef{列上 18:34} 因为如果我们承认这一点，并且脱离行事者的意图，单就行为本身来考察，那么如果有人愿意，他就可以判\Person{亚巴郎}犯了杀子之罪 \BibleRef{创 22:3}，并控告他的孙子和后裔 \BibleRef{出 11:2} 犯了邪恶和诡诈之罪。因为前者获得了长子的名分，后者将埃及人的财富转移到了以色列子民的营地。但事实并非如此：抛弃这大胆的想法！因为我们不仅免除了他们的罪责，而且因这些事而赞美他们，因为连天主也曾为此赞扬过他们。因为那不正当地使用这种手段的人，才应恰当地被称为欺骗者，而不是为了有益目的这样做的人。而且常常有必要欺骗，并通过这种手段行最大的善事，而直来直去的人却对他没有欺骗的人造成了巨大的伤害。\par

\SourceRecord{Translated by W.R.W. Stephens.；Nicene and Post-Nicene Fathers, First Series；Vol. 9.；Edited by Philip Schaff.；Buffalo, NY: Christian Literature Publishing Co.,；1889.；<http://www.newadvent.org/fathers/19221.htm>.}

% structure: p0001 source=h1 target=section reason=page-title

\section{论司铎职（卷二）}

1. 那么，为善的目的而使用欺骗是可能的，或者更确切地说，在这种情况下不应称之为欺骗，而是一种值得赞赏的妥善处置，这一点本可更详细地证明；但既然已说的足以证明，若再延长这个话题，就会令人厌烦且冗长。现在，需要你来证明我是否没有利用这种技巧为你谋取益处。\par

\Person{巴西略}: 那么我从这妥善处置、或明智策略、或你乐意称之为的任何东西中得到了什么益处，以致让我相信没有被你欺骗呢？\par

\Person{金口若望}: 请问，还有什么益处能比被人看见做那些基督亲口宣称是爱祂的明证的事更大呢？\BibleRef{若 21:15} 因为祂对宗徒之长说：\BibleQuote{伯多禄，你爱我吗？}当他承认时，主接着说：\BibleQuote{如果你爱我，就牧放我的羊群。}师傅问门徒是否被爱，不是为了获取信息（那洞察人心者怎能不知？），而是为了教导我们祂对牧养这些羊群有多么大的关注。这一点既已明了，同样显而易见的是，凡为这些基督如此珍视的羊群劳苦的人，将获得大而不可言喻的赏报。因为当我们看见有人照顾我们的家眷或羊群时，便将他对此的热忱视为爱我们的标记：然而这一切都可以用金钱买到——那么祂将用多么伟大的恩赐来回报那些牧放祂所买来的羊群的人呢？这羊群并非用金钱或类似之物买来，而是用自己的死亡，以祂的血作为羊群的代价。因此，当门徒说：\BibleQuote{主，祢知道我爱祢，}并请所爱者亲自为他的爱作证时，救主并没有停在那里，而是加上了爱的记号。因为祂那时并非要显示伯多禄多么爱祂，而是要显示祂自己多么爱祂的教会，并愿意教导伯多禄和我们众人，我们也当为此付出同样大的热忱。因为天主为何没有怜惜自己的独生子，反而将祂交出，尽管祂只有这独子？是为了使那些与祂为敌的人与祂和好，并使他们成为自己的子民。祂流血的目的是什么？是为了赢得这些祂托付给伯多禄及其继承者的羊群。于是基督自然说：\BibleQuote{那么谁是那忠信而聪明的仆人，主人将派他管理自己的家仆呢？}再者，这些话看似疑惑，但说话者并非出于疑惑而说，正如祂问伯多禄是否爱祂，并非需要知道门徒的情感，而是出于彰显自己爱情的极度深广：同样，当祂说：\BibleQuote{那么谁是那忠信而聪明的仆人？}祂说这话并非不知谁是忠信聪明的，而是渴望表明这种品格的稀有和这职位的伟大。无论如何，请看赏报何其大——祂说：\BibleQuote{祂要指派他管理他的一切财产。}\BibleRef{玛 24:47}\par

2. 那么，你还会坚持说你没有被正当地欺骗吗？因为你将管理属于天主的事物，而且正在做那事——当\Person{伯多禄}做了这事时，主说他将能够超越其他宗徒，因为祂的话说：\BibleQuote{伯多禄，你爱我比这些更深吗？}然而祂本可以对他说：\Quote{如果你爱我，就操练禁食、睡在地上、长时间守夜、保护受屈者、作孤儿的父亲、代替丈夫照顾他们的母亲。}但实际上，撇开所有这些，祂说什么？\BibleQuote{牧放我的羊。}因为我已提过的那些事，即使是许多在权下的人，无论男女，也容易做到；但当一个人需要管理教会并被托付照顾这么多灵魂时，整个女性性别在面对如此重大的任务时必须退后，大多数男性也是如此；我们必须提出那些远超过其他人的人，他们在精神上的卓越如此超越，如同\Person{撒乌耳}在身体上超过全体希伯来民族一样，甚至更甚。\BibleRef{撒上10:23}因为在此情况下，不要让我以肩膀的高度作为探究的标准；而要让他与他的羊群之间的差别，如同理性的人与无理性的受造物之间的差别一样大，甚至更大，因为风险涉及更加重要的事物。他若失去羊群，或由于狼的蹂躏，或盗贼的攻击，或瘟疫，或其他灾难降临，或许能从羊群的主人那里得到一些宽恕；即使主人要求赔偿，处罚也只是金钱的事。但受托管理人类——基督的理性羊群——的人，首先要为失去羊群而受罚，这超过物质事物而触及他自己的生命；其次，他必须进行一场更伟大、更困难的斗争。因为他不是与狼搏斗，也不是惧怕盗贼，也不是考虑如何使羊群免于瘟疫。那么，他要与谁作战？与谁角力？请听\Person{圣保禄}的话：\BibleQuote{我们不是与血肉之敌搏斗，而是与率领者、与掌权者、与这个黑暗世界的统治者、与天界中的邪恶势力搏斗。}\BibleRef{弗6:12}你看到可怕的多重敌人了吗？看见他们凶猛的队伍了吗？不是穿着铁甲，而是具有一种本性，本身就相当于全套盔甲。你想再看另一支队伍，严厉而残酷，围攻这羊群吗？你也要从同一观察点看见它。因为那位向我们讲述其他敌人的人，也向我们指出这些敌人，在某个地方这样说：\BibleQuote{肉体的行为是明显的，即：淫乱、奸淫、不洁、放荡、崇拜偶像、施行邪术、仇恨、争吵、嫉妒、愤怒、争执、诽谤、窃窃私语、自高自大、混乱，}\BibleRef{格后12:20}还有许多；因为他没有列全，而是让我们从这些中明白其余的。此外，对于无理性的羊群的牧人，那些想要毁灭羊群的人，当他们看见守护者逃跑时，就不再与他作战，满足于抢夺牲畜。但在这种情况下，即使他们捕获了整个羊群，他们也不会放过牧人，反而更加攻击他，更加胆大妄为，直到要么推翻他，要么自己被击败。此外，羊群的不幸是明显的，无论是饥荒、瘟疫、创伤，还是其他任何使之痛苦的事，这或许对那些在这方面受苦者获得缓解不无帮助。还有一个更大的事实有助于解除这种软弱。那是什么？牧人有权柄强制羊群接受治疗，即使它们不自愿。因为当需要灼烧或切割时，容易将它们绑住，并将它们关在圈内很长一段时间，只要有益处，并给它们另一种食物代替另一种，切断它们的水源，以及牧人认为有利于它们健康的所有其他事情，他们都能很容易地做到。\par

3. 但就人的软弱而言，首先人不容易辨别它们，因为没有人\BibleQuote{知道人的事，除非人里面的那人的心神。} \BibleRef{格前 2:11} 那么，一个人若不晓得病症的性质，又如何能对症下药呢？常常连他自己若染上此病也无法理解。即使疾病变得明显，也给他带来更多麻烦：因为不可能像牧人对付羊那样，以同样的权威医治所有人。因为在这种情况下，也需要捆绑、限制食物，并使用烙铁或刀；但治疗的接受取决于病人的意愿，而非施治者的意愿。那位奇妙的人（\Person{圣保禄}）也领悟到这一点，他对格林多人说：\BibleQuote{并不是我们管制你们的信德，而是我们协助你们的喜乐。} \BibleRef{格后 1:24} 因为基督徒比其他任何人都更不被允许强制纠正犯罪者的过失。世俗的审判官确实在依法捕获罪犯后，显示其权威很大，并阻止他们违心地追随自己的私欲；但在我们这里，作恶者必须通过说服，而非通过强迫来改善。因为这样的限制罪人的权威并没有通过法律赐予我们，即使赐予了，我们也没有施展权力的余地，因为天主奖赏那些自愿而非被迫远离邪恶的人。因此，需要极大的技巧，好使病人自愿接受医生所开的治疗，不仅如此，还要使他们为痊愈而感恩。因为如果有人被捆绑时变得倔强（这在他能力之内），他会使危害加重；而如果他不留意那些如刀割般的话语，他会因这种藐视而再添新伤，医治的意图反而成了更严重混乱的起因。因为不可能强迫任何人违背其意愿来医治他。\par

4. 那么该怎么办呢？因为如果对一个需要严格动刀的人处理得太温柔，没有深深切入需要治疗之处，你除去了一部分病灶，却留下了另一部分；另一方面，如果你毫不留情地做了必要的切口，病人被痛苦逼得绝望，常常会一下子把一切——药物和绷带——都扔掉，并一头栽下去，\Quote{折断轭，挣断绳索。} 我能说出许多人因被追究应有的惩罚而陷入极端邪恶的事。因为我们在施行惩罚时，不应该仅仅按照罪过的程度来量刑，而更应考虑罪人的心态，免得我们想修补裂口，却使裂口更甚，热心恢复倒塌的，反而使废墟更大。因为软弱和粗心的人，大多沉溺于世俗享乐，并因出身和地位而骄傲，如果温和地、逐渐地引导他们悔改自己的错误，至少能够部分地、即便不是完全地，从他们所占有的邪恶中得到解脱；但如果有人一下子施以纪律，就会夺去他们这微小的改正机会。因为一旦灵魂被迫抛弃羞耻，它就陷入麻木状态，既不听从温和的话，也不屈服于威胁，也不接受感激，反而比先知斥责的那个城市更坏，他说：\BibleQuote{你却有娼妓的面孔，拒绝在人前羞愧。} \BibleRef{耶 3:3} 因此，牧人需要极大的谨慎，并要有无数的眼睛从各方面观察灵魂的习惯。因为正如许多人被提升到骄傲，然后因无法忍受严厉的疗法而沉入对自己救恩的绝望，也有一些人由于未受到与罪相应的惩罚而陷入疏忽，变得更坏，并被推向更大的罪。因此司铎应当不遗漏这些事中的任何一件，而是在彻底调查所有情况后，针对每个案件恰当地使用疗法，免得他的热忱徒劳无功。不只在这一点上，而且在弥合教会分离的肢体上，也能看到他有许多工作。因为牧人的羊群跟随他，无论他领到哪里；如果有人偏离了直路，离开良好的牧场，在贫瘠或崎岖的地方吃草，一声响亮的呼喊就足以集合它们，并把那些离群的羊带回到羊栈；但如果是一个人偏离了正确的信仰，就需要极大的努力、恒心和耐心；因为他不能被强力拖回，也不能被恐惧约束，而必须通过说服领回那最初偏离的真理。因此，牧人应当有高尚的精神，不致气馁或对离群者的得救绝望，而应不断地对自己说：\BibleQuote{或许天主会赐给他们悔改，得以认识真理，并摆脱魔鬼的罗网。} \BibleRef{弟后 2:25} 因此主在对门徒们说话时说：\BibleQuote{那么，谁是忠信而聪明的仆人？} \BibleRef{玛 24:45} 因为那操练自己的人只关心自己的益处，而牧养职务的益处却延伸到全体人民。而且，向穷人施舍金钱或以其他方式帮助受压迫的人，也在某种程度上惠及邻人，但比司铎的益处小得多，正如身体低于灵魂。所以主正确地指出，对羊群的热忱是爱祂的标志。\par

\Person{圣巴西略}：但你本人——难道你不爱基督吗？\par

\Person{金口圣若望}：是的，我爱祂，永不停止爱祂；但我恐怕会触怒我所爱的祂。\par

巴西略：但有什么谜语比这件事更难解呢？基督命令爱祂的人牧放祂的羊，而你却说因为爱那颁下这道命令的主，所以你拒绝牧放它们？\par

金口若望：我说的话并非谜语，而是非常明白简单的。因为如果我有充分的资格管理这职务，正如基督所期望的，而我又回避它，我的说法或许可被质疑；但既然我精神的软弱使我对这职务一无所用，为什么我的话该被怀疑呢？因为我害怕，如果我将那原本状况良好、喂养充足的羊群接手，却因我的不熟练而浪费它们，我会激起那如此爱羊群以至于为它们的救恩和赎价而舍己的天主对我的义怒。\par

巴西略：你在开玩笑吧！因为如果你是认真的，我不知道除了用你试图驱散我沮丧的这些话，你还能如何证明我应受这样的悲伤。我原本就知道你欺骗并出卖了我，但现在当你着手为我对你的指控辩白时，我更清楚地看到和理解你使我陷入何等大的灾祸。因为如果你因为意识到自己的精神不足以承担这重任而退出这职务，那么即使我碰巧非常渴望这职务，我也应该在你之前被解救出来，更不用说我把这些事的全部决定权都交给了你；但事实是，你只考虑自己的利益而忽视了我。你若完全忽视我才好呢！那样我会非常满意；可是你谋划着方便那些想要抓我的人将我捕获。因为你不能以舆论欺骗了你、使你对我产生巨大而奇妙的想法为托词。因为我不是你们那些了不起的杰出人物之一，就算我是，你也不该将舆论置于真理之上。因为我若从未允许你享受我的陪伴，你或许似乎有合理的借口根据公众传闻来投票；但若没有人像我一样如此透彻地了解我的事情，你若比我的父母和养育我的人更清楚我的品性，那么你能用什么足够令人信服的论证来说服听众，说你并非故意把我推入这危险之中？说，我该如何回应那些指控你的人呢？\par

金口若望：不！除非我先解决了那些只关乎你自身的问题，我不会继续讨论这些问题，即使你要求我千百次处理这些指控。你说无知会让我获得宽恕，若我对你一无所知而把你带入现在的处境，我会免于一切指控；但因为我不是在无知中出卖你，而是熟悉你的事情，我便失去了所有合理的借口和托辞。但我说恰恰相反：因为在这类事上，需要仔细审查，凡要举荐某人胜任司铎职务的人，不应只满足于公众传闻，而应首先、并在一切之前亲自调查那人的品性。因为当蒙福的保禄说：\Quote{他也必须在那些外人中有好名声}，\BibleRef{弟前 3:7} 他并未免除精确而严格的审查，也未将这种见证置于这类事所需的考察之上。因为在作了许多前导论述之后，他提到了这个附加的见证，证明人不可只满足于它来进行这种选举，而应与其余的一同考虑。因为公众传闻常常说谎；但当仔细考察在先时，就不再需要担心它了。为此，他在其他证据之后列出了来自外人的证据。因为他不是简单地说：\Quote{他必须有好的名声}，而是加上了：\Quote{来自那些外人}，想表明在来自外人的名声之前，他必须被仔细审查。因此，既然我自己比你父母更了解你的事情，正如你自己也承认的，我或许该被免于一切责备。\par

巴西略：不，这正是为什么如果有人要控告你，你无法逃脱。难道你不记得曾听我说过，并经常从我的实际行为中了解到我性格的软弱吗？你不是一直因我的懦弱而嘲弄我，因为我那么容易因日常忧虑而气馁吗？\par

5. 金口若望：我确实记得常常听你说过这样的话；我不否认。但如果我曾嘲弄你，那是出于玩笑而非认真的实话。不过，我现在不争论这些事，我要求你同样容忍，因为我想提及你拥有的一些优点。因为如果你试图证明我说的是假的，我不会放过你，而会证明你这样说更多出于自我贬抑而非为了真理，我将只用你自己的言语和行为来证明我论断的真实性。现在，我首先想问你的问题是：你知道爱的力量有多大吗？因为省略了宗徒们将要行的所有奇迹，基督说：\Quote{如果你们彼此相爱，众人因此就可认出你们是我的门徒了}，\BibleRef{若 13:35}；保禄说爱是法律的圆满，\BibleRef{罗 13:10}，没有爱任何神恩都没有益处。那么，这精选的善、基督门徒的标志、高于所有其他恩赐的恩赐，我看出它深深植根于你的灵魂，并结出丰硕的果实。\par

巴西略：我承认这对我确实是一件深切关注的事，并且我极其恳切地努力遵守这条诫命，但我甚至没有做到一半，即使你自己也会为我作证，如果你不再出于偏袒说话，而只是尊重真理。\par

6. \Person{金口圣若望}：那么，我现在就来引用证据，并履行我先前的威胁，证明你是想要贬低自己，而非说出真相。但我要提一件刚刚才发生的事，免得有人怀疑我因年代久远而试图模糊真相，因为遗忘会使人无法对我为了取悦而说的话提出异议。有一位我们亲密的朋友，被人诬告侮辱和愚妄，陷入了极大的危险之中，那时你挺身而出，投身于危险之中，尽管没有人召唤你，也没有那位即将陷入危险的人向你求助。事实就是如此；为了用你自己的话来定你的罪，我要提醒你所说的话：当有些人不赞成这种热忱，而另一些人称赞和钦佩它时，你对那些指责你的人说：\Quote{我怎能不这样做呢？因为我不知道除了在必要时刻为拯救任何身陷危险的朋友而舍弃生命之外，还有什么别的爱的方式：}这样，你用不同的言辞，却以同样的意义，重复了\Person{基督}在向门徒们阐述完美之爱的定义时所说的话。祂说：\BibleQuote{人若是为朋友舍弃生命，再没有比这更大的爱了。}如果不可能找到比这更大的爱，那么你已经达到了它的极限，并以言行加冕了巅峰。这就是我出卖你的原因，这就是我设下那计谋的原因。现在我可曾说服你，我并非出于任何恶意，也不是因为想要把你推入危险，而是出于对你未来益处的确信，才把你拖入这条道路？\par

\Person{巴西略}：那么你认为爱足够用来纠正同伴吗？\par

\Person{金口圣若望}：当然，这在很大程度上会有助于此。但如果你希望我也拿出你实践智慧的证据，我就来做，并证明你的理解力超过了你的仁爱。\par

听到这些话，他满脸通红地说：现在不要再谈我的品性了吧；因为我最初要求解释的并不是这个；但如果你对那些外人有什么合理的回答，我很乐意听听你要说什么。所以，放弃这无谓的争论吧，请告诉我，我该如何向那些尊重你的人以及那些因他们认为自己受到侮辱而伤心的人辩护。\par

7. \Person{金口圣若望}：这正是我最终要讲到的点，因为我对你的解释已经完成，我将很容易转向辩护的这一部分。那么这些人的控告是什么？他们的指控是什么？他们说我侮辱了他们，严重地亏待了他们，因为我没有接受他们想赐予我的荣誉。首先，我要说，不应考虑对人的侮辱，因为如果我去尊敬他们，我就不得不冒犯\Person{天主}。我要对那些不快的人说，为这些事情生气是不安全的，反而会给他们带来很大的伤害。因为我认为，那些倚靠\Person{天主}、只仰望\Person{天主}的人，应当如此虔诚，即使他们被羞辱一千次，也不应将其视为侮辱。但我甚至从未想过敢做这样的事，从我将要说的话中可以明显看出。因为如果我确实是出于傲慢和虚荣（就像你常说的有人诽谤那样）而同意了我的控告者，那我就是最邪恶的人之一，轻慢了伟大而优秀的人，况且他们还是我的恩人。如果人们因得罪从未得罪他们的人而应受惩罚，那么对于自愿优先尊重我们的人，我们该如何尊敬他们呢？因为没有人能说他们是在回报我从他们那里得到的任何大小恩惠。那么，如果一个人以相反的方式回报他们，该受多大的惩罚呢？但如果我从未想过这样的事，而是出于完全不同的意图拒绝了这沉重的负担，那么为什么他们拒绝原谅我（即使他们不同意赞许），却指责我自私地保全了自己的灵魂？因为远非侮辱了那些人，我甚至要说，我通过拒绝反而尊重了他们。\par

不要对我的话的悖论性质感到惊讶，因为我马上就会给出一个快速的解答。\par

8. 因为我若接受了职务，且不说所有人，就是那些喜欢说坏话的人，也可能对我这个被选者以及他们这些选举者怀疑并说出许多事：例如，他们看重财富，羡慕显赫地位，或受阿谀奉承而提拔我获得此尊严；其实我无法确定是否有人会怀疑他们被金钱收买。此外，他们会说：\Quote{基督召叫渔夫、制帐篷者和税吏来获得此尊严，而这些人却拒绝那些靠每日劳作谋生的人；但若有任何人投身于世俗学问，并在闲散中长大，他们就接纳并赞赏。请问，为什么他们忽略了那些为教会服务而经历无数辛劳的人，却突然将这个从未经历过此类辛劳、却将全部青春耗费在世俗学问的虚妄研究中的人拖进这个尊严？} 若我接受了职务，他们会说这些甚至更多；但现在情况并非如此。因为他们所有毁谤的借口都被切断了，他们既不能指责我阿谀奉承，也不能指责他人收受贿赂，除非有人像疯了一样。因为一个用奉承和花钱来获得尊严的人，在可以得到它时，却把它让给了别人，这怎么可能？这就像一个人为土地付出了很多劳力，希望田地丰收，酒榨流溢，经过无数辛劳和巨大花费后，却在收获庄稼和采集葡萄的时候把果实让给别人。你看，虽然所说的可能远离真相，但那些想要诽谤选举者的人那时会有借口声称选择是不公正且未经考虑的。但如今我阻止了他们张口，甚至在此事上说不出一句话。那么，最初他们的评论就是如此甚至更多。但若我承担了牧职，我无法日复一日地为自己辩护，对抗控告者，即使我一切做得无懈可击，更不用说由于我的年轻和缺乏经验而必定会犯的许多错误。但如今，我也使选举者免于这种指控，而在另一种情况下，我会使他们陷入无数的指责。因为世人会怎么说？\Quote{他们将如此重大且重要的事务托付给轻率的年轻人，他们玷污了天主的羊群，基督徒的事务成了笑柄。} 但现在\Quote{一切不义都要闭口无言。} 因为虽然他们可能因你们而说这些话，你们将迅速以行动教导他们：理解不因年龄衡量，白发不是检验年长者的标准——年轻人不应被绝对排除在牧职之外，只有初信者才被排除：两者之间区别很大。\newline{}\par

\SourceRecord{Translated by W.R.W. Stephens.；Nicene and Post-Nicene Fathers, First Series；Vol. 9.；Edited by Philip Schaff.；Buffalo, NY: Christian Literature Publishing Co.,；1889.；<http://www.newadvent.org/fathers/19222.htm>.}

% structure: p0001 source=h1 target=section reason=page-title

\section{论司铎职（第三卷）}

1. \Person{金口若望}：关于那些侮辱我尊荣之人的事，我先前所述足以证明，我逃避此职并非有意羞辱他们；但我现在将尽己所能，使事情显明，我并非因任何骄傲而自大。因为倘若有人将将军或王国的抉择交予我，而我当时已下定决心，任何人或许自然怀疑我有此过错，甚或所有人都会判定我并非骄傲，而是愚昧。但当司铎职授予我时，它远超王位，如同精神与肉体之别，难道还有人敢指控我轻视吗？拒绝小事之人被指为愚蠢，而在极其重要之事上如此行的人，却免于精神错乱的指控，反而被指控骄傲，这岂非荒谬？这如同有人指控一个人轻视牛群、拒绝为牧人，不是说他骄傲，而是说他疯狂；但若有人拒绝世界帝国及地上所有军队的统帅权，却说此人并非疯狂，而是骄傲膨胀。但情况绝非如此；说这些话的人伤害他们自身远胜于伤害我。因为仅仅想象人性可能轻视此尊位，便足以证明那些提出此指控的人对该职务的评价。因为若他们不视此为平常无足轻重之事，这种怀疑绝不会进入他们的头脑。为何从未有人敢于对天使的尊位抱有此种怀疑，说人性不愿向往天使的本性是因为骄傲？这是因为我们对于那些能力怀有伟大想象，这不容我们相信人能想像出比那尊荣更伟大的事物。因此，那些指控我骄傲的人，倒更有理由被我控告骄傲。因为若非他们先前已贬低此事为无足轻重，他们绝不会对他人产生这种怀疑。但若他们说我是为了光荣而行此事，他们将被证实是公然自相矛盾，自投罗网；因为我不晓得他们能找到何种论据（若他们想使我免于虚荣的指控）比这更合适。\par

2. 因为若这种渴望曾进入我的思想，我本该接受而非逃避此职。为何？因为它会给我带来极大光荣。因为像我这般年纪的人，如此近期放弃世俗追求，却突然被众人认为值得如此钦佩，以至在那些终生从事此类劳作的人之前获得荣耀，并获得比他们所有人更多的票数，本可说服所有人预期我会有伟大奇妙之事。但事实上，教会中大部分甚至不知我名：因此，我拒绝此职甚至不会为众人所知，仅限于少数人，而我甚至不确定这些人中是否所有人都确切知晓；但很可能他们中许多人或者以为我根本未被选上，或者以为我是在选举后被拒绝，被认为不合格，而非我自愿逃避此职。\par

3. \Person{巴西略}：但那些知晓真相的人将会惊讶。\par

\Person{金口若望}：看哪！这些人就是，按你所说，错误地指控我虚荣和骄傲的人。那么我该从哪里期待赞美？从多数人？他们不知道实际情况。从少数人？这里事情又被扭曲而于我不利。因为你此刻来此的唯一理由，就是应当学习该给他们什么答复。那么，鉴于这些事，我现在该确定地说什么？因为稍等片刻，你将清楚察觉，即便所有的人都知晓真相，他们也不应当因骄傲和贪求光荣而谴责我。此外还有另一层考虑：不仅那些做出这种指控的人（若有任何这样的人——我本人并不相信），而且那些对他人抱此怀疑的人，都将陷入不小的危险。\par

4. 因为司铎职务确实在地上执行，但位列于天上的规章之中；这是非常自然的：因为并非人、天使、总领天使或任何其他受造之力，而是护慰者祂自己建立了这一圣召，并说服仍在肉身中的人代表天使的职务。因此，祝圣的司铎应当纯洁，仿佛他本人站在天上，处于那些能力之中。确实，在恩宠的时期之前所使用的事物，如铃铛、石榴、胸牌和厄弗得上的宝石、腰带、头巾、长袍、金牌、至圣所、里面的深深寂静，都是可畏且极其庄重的。但是，倘若有人审察属于恩宠的时期的事物，他会发现，它们虽小，却同样可畏且充满敬畏，并且关于法律所说的这话在此也同样真实：\BibleQuote{那先前有光荣的，因了这更超越的光荣，已算不得光荣了。} \BibleRef{格后 3:10} 因为当你看见主被祭献，放置在祭台上，司铎站在祭品前祈祷，所有参礼者都被那宝血染红，你还会认为你仍在人间，站立在地上吗？难道你不是相反地立刻被提升到天堂，从灵魂中驱除一切肉性的思想，以脱去形体的精神和纯理性默观天上的事物吗？哦！何等奇事！天主对人的爱何等伟大！祂与圣父同坐在高天，在那时刻被众人所持有，并将自己交给那些愿意拥抱和把握祂的人。而这一切都是通过信德的眼睛做到的！这些事在你看来是可轻视的，或使人可能因它们而骄傲吗？\par

你岂不愿从另一个奇迹中学习这职务的卓绝神圣么？试想\Person{厄里亚}与环绕他的庞大群众，祭品置于石砌的祭台上，其余百姓都深深静默，唯有先知独自信祷：然后火焰突然从天降于祭品上——这些是奇妙可畏的事，充满恐怖。现在请由此场景转到今日所举行的礼节；它们不仅奇妙可睹，更超越恐怖。那里站着司铎，并非从天上引下烈火，而是引下圣神；他作长久的恳求，并非祈求从高天降下的火焰焚烧供物，而是祈求恩宠降临于祭品，藉以光照众人的灵魂，使他们比经火炼净的银子更加璀璨。谁能轻蔑这至可敬畏的奥迹，除非他完全疯狂失去理智？或者你不知道，若非天主恩宠的大力援助，没有一个人的灵魂能承受那祭火，而全都要被彻底焚毁么？\par

5. 因为如果有人思量，一个人，有血肉之躯，得以亲近那有福而纯洁的本性，是何等大事，那么他将清楚看见圣神的恩宠赐予了司铎何等大的荣耀；因为藉着他们的服务，这些礼节得以举行，还有其他与此不相上下的礼节，关系我们的尊严与救恩。因为那些居住在地上并以此为居处的人，被托付管理天上之事，并获得了天主不曾赐给天使或总领天使的权力。因为从未对他们说过：\Quote{凡你在地上束缚的，在天上也要被束缚；凡你在地上释放的，在天上也要被释放。}\BibleRef{玛 18:18} 那些在地上掌权的，确实有束缚的权力，但仅限于身体；而这里的束缚却抓住灵魂并穿透诸天；司铎在此世所做的一切，天主在上天予以认可，主人确认其仆人的判决。因为当祂说：\Quote{你们赦免谁的罪，就给谁赦免；你们保留谁的罪，就给谁保留？}\BibleRef{若 20:23} 还有什么权柄比这更大？\Quote{父把一切审判交给了圣子？}\BibleRef{若 5:22} 但我看见这一切都由圣子交给了这些人。因为他们被带领到如此尊严，仿佛已被提升到天堂，超越了人性，并摆脱了我们所易受的情欲。再者，如果一位国王将这尊荣赐予他的一位臣民，授权他随心所欲地下狱或释放，那么此人便成为众人羡慕尊敬的对象；而那从天主接受了比这更大权柄的人（因为天比地更宝贵，灵魂比身体更宝贵），有些人竟认为他接受了如此小的尊荣，以至于他们竟能想象那些受托这些事的人中会有人轻视这恩赐。愿这样的疯狂离开！因为轻视如此伟大的尊严，就是明显的疯狂；没有这尊严，就不可能获得我们自己的救恩，也不可能获得那应许给我们的美物。因为若没有人能进入天国，除非他藉着水和圣神重生；并且那不吃主的肉、不喝祂的血的，必被排除于永生之外；而这一切都唯有藉着那些圣洁的手——我是指司铎的手——才能完成；那么没有这些，谁能逃脱地狱之火，或赢得那为得胜者保留的冠冕呢？\par

6. 这些人确实受托负责属灵的产痛和藉着圣洗而来的重生：藉着他们，我们穿上基督，与天主子同葬，成为那蒙福之首的肢体。因此，他们不仅比统治者和君王更应受我们敬畏，而且比父母更应受尊敬；因为父母生于血和肉情，而他们却是我们由天主而生的根源，即那蒙福的重生，它是真正的自由和按恩宠而来的嗣子身份。犹太司铎有权释放身体上的癞病，或者更确切地说，并非释放，而只是检验那些已被释放的人，你们知道当时司铎的职务是多么被争夺。但我们的司铎已获授权处理，不是身体上的癞病，而是属灵的不洁——不是检验后宣布去除，而是实际上绝对地除去它。因此，藐视这些司铎的人，比达堂及其同伙更该受诅咒，应受更严厉的惩罚。因为后者虽觊觎不属于自己的尊位，但仍对之有极好的看法，并藉着追求它时极大的热忱证明了这一点；而这些人，当职务已更好地调整并得到如此巨大的发展时，却表现出一种超过前者的鲁莽，尽管是以相反的方式表现出来。因为觊觎不属于自己的尊荣，与藐视如此巨大的利益，所含的轻蔑并不等同，后者超过前者正如轻视有别于钦佩。那么，有什么灵魂如此卑下，竟藐视如此巨大的利益呢？我认为一个也没有，除非是受某种魔鬼冲动支配的灵魂。因为我要再次回到我出发的那一点：天主赐给司铎一种大于我们生身父母的权能，不仅在惩戒方面，而且在施惠方面。两者确实有多大差别，就像现世与来生一样。因为我们的生身父母只为我们生于此世，而后者却为那将来之世。前者无法使他们的后代免于死亡，也无法抵御疾病的侵袭；但后者却常常拯救有病的灵魂，或即将灭亡的灵魂，为一些人获得较轻的惩罚，并阻止另一些人完全堕落，不仅通过教导和劝诫，也通过祈祷所获得的援助。因为不仅在重生之时，而且在之后，他们也有权赦罪。经上说：\BibleQuote{你们中间有患病的吗？他该请教会的长老们来，让他们为他祈祷，因主的名给他傅油。出于信德的祈祷，必救那病人，主必使他起来；如果他犯了罪，也必得赦免。} \BibleRef{雅 5:14}再者，我们的生身父母，若他们的子女与世上任何位高权重的人发生冲突，便无法帮助他们；但司铎却曾使那多次被他们激怒的天主和好——不是使统治者和君王和好，而是使天主自己和好如初。\par

那么！此后还有谁敢谴责我傲慢呢？就我而言，在说了这些之后，我想如此宗教性的敬畏将占据听众的灵魂，使他们不再谴责那些因傲慢和轻率而回避职务的人，反而谴责那些自愿上前、急切为自己求得这尊位的人。因为那些受托管理城市的人，若偶然缺乏谨慎和警醒，有时会毁灭城市并同时毁灭自己；那么，你们认为，那负责装饰基督新娘的人，为避免犯罪，需要多少自身与来自上天的力量呢？\par

7. 没有人比保禄更爱基督：没有人表现出更大的热忱，没有人被认为配得上更多的恩宠：然而，尽管有这些巨大的优势，他仍然为这治理及那些被他治理的人感到恐惧战栗。他说：\\BibleQuote{我怕你们的心意，或许像蛇用诡诈诱惑了厄娃一样，以致腐化了那在基督内的纯一。} \\BibleRef{格后 11:3} 又说：\\BibleQuote{我在你们那里时，又恐惧，又甚战栗；} \\BibleRef{格前 2:3} 而这人曾被提升到第三层天上，并分享了天主不可言传的奥秘\\BibleRef{格后 12:4}，他自从成为信徒后，每天历尽如同日子一样多的死亡——而且，这人还不愿使用基督赐予他的权柄，免得他的任何皈依者跌倒。那么，若他超越了天主的命令，从不寻求自己的益处，只求他所治理者的益处，却仍因考虑这治理的重大而常充满恐惧，我们的情况又将如何？我们在许多方面寻求自己的益处，不仅未能超越基督的命令，反而大部分时候违犯它们。他说：\\BibleQuote{有谁软弱，我不软弱呢？有谁跌倒，我不焦急呢？} \\BibleRef{格后 11:29} 司铎应当是这样的人，或者，不如说，不仅是这样的：因为与我即将说的相比，这些是小事，微不足道。这是什么？他说：\\BibleQuote{为了我的弟兄，我血统上的同胞，我宁愿自己受诅咒，与基督隔绝。} \\BibleRef{罗 9:3} 若有人能说出这样的话，若有人拥有能达到如此祈祷的灵魂，那他若逃避职分尚可责备；但若有人像我这样缺乏如此卓越的品德，那么他不是逃避，而是接受这职分才该被憎恨。因为如果当前的事务是选举一位军事将领，而有权授予这荣誉的人拉来一个铜匠、鞋匠或某个工匠，将军队托付给他，我不会称赞那可怜的人，如果他没有逃跑并尽力避免陷入如此明显的麻烦。倘若仅仅拥有牧者的名号、随意处理工作便足够，且这毫无危险，那么谁想指责我虚荣，就尽管指责吧；但若承担这照料的人需要许多智慧，在智慧之前更需要来自天主的巨大恩宠、正直的品行、纯洁的生活和超性的德行，那么请别剥夺我的宽恕，如果我不愿无谓地枉自丧亡。\par

此外，若有人负责一艘满载划桨手、载有昂贵货物的大商船，命我掌舵，并要我穿越\\Place{爱琴海}或\\Place{第勒尼安海}，我会立刻拒绝这提议；若有人问我为什么？我会说：\\Quote{免得我使船沉没。} 好吧，当损失涉及物质财富，危险仅止于肉身的死亡时，无人会责备那些极度谨慎的人；但若沉船者注定不是坠入海洋，而是坠入永火深渊，且等待他们的死亡不是使灵魂与肉身分离，而是连同此世一并打入永罚，我岂会因未一头栽入如此巨大的灾祸而惹你发怒和憎恨呢？\par

8. 请不要这样，我恳求并央求你们。我认识自己的灵魂，它是多么软弱无力：我知道这职务的重大，以及工作的巨大困难；因为司铎的灵魂受到比惊扰海洋的狂风更猛烈的波涛折磨。\par

9. 首先，最可怕的是虚荣的暗礁，比神话编造者所讲述的\\OriginalTerm{塞壬}{Sirens}的暗礁更加危险：许多人能驶过那暗礁而毫发无损；但这暗礁对我如此危险，以致如今，当没有任何必要驱使我进入那深渊时，我仍无法避开这陷阱；但若有人将这职责托付给我，那就如同绑住我的双手，将我交给栖息在那暗礁上的野兽，让它们日复一日地撕裂我。你问这些野兽是什么？它们是愤怒、沮丧、嫉妒、纷争、诽谤、控诉、谎言、伪善、阴谋、对无辜者的怒气、对同工不当行为的快感、对他们的兴旺的忧愁、对赞美的爱、对荣誉的渴求（这最会将人的灵魂推入毁灭）、讨好人的教义、卑贱的奉承、下作的谄媚、对穷人的蔑视、对富人的巴结、无意义且有害的荣誉、对给予者和接受者都危险的恩惠、只配最下贱奴隶的卑污恐惧、直言的无制、对谦卑的假装、对真理的放逐、对谴责与规劝的压制，或更确切地说，对穷人过度加以谴责，而对有权势者却无人敢开口。\par

所有这些野兽，以及更多，都孕育在我所说的那暗礁上，一旦被它们捕获，便不可避免地陷入如此深的奴役，以致为了取悦妇女，她们常做许多不便于提及的事情。神圣的法律已将妇女排除在职务之外，但她们却试图闯入其中；既然她们自己无法成就什么，她们便通过他人做一切事；并且她们已拥有如此大的权力，以至于能按自己的意愿任命或驱逐司铎；事实上，事情被颠倒了，如谚语所说：\\Quote{被管的人管带领袖：} 但愿这是男人而非女人做的事，因为女人没有被授予教导的职务。为何我说教导？因为蒙福的保禄甚至不允许她们在教会中说话。但我曾听说，她们已获得如此大的言论自由，甚至能责备教会的主教们，比主人责骂自己的仆人更加严厉。\par

10. 并且请任何人不要认为我把所有人都归入上述指控：因为有些人，是的，许多人，优于这些缠累，其数量超过被这些缠累所捕获的人。我也不愿让司铎职为这些邪恶负责：愿这样的疯狂远离我。因为明白事理的人不说刀剑应对谋杀负责，也不说酒应对酗酒负责，也不说力量应对暴行负责，也不说勇力应对鲁莽负责，而是归咎于那些滥用天主所赐恩惠的人，并相应地惩罚他们。确实，至少司铎职可以公正地控告我们，如果我们没有正确履行它。因为它本身并不是上述邪恶的原因，而是我们，在我们能力范围内，用如此多的污秽玷污了它，将其托付给那些轻易接受所提供之物的平庸之人，他们事先没有获得对自己灵魂的认识，也没有考虑职务的严重性，当他们进入工作时，因缺乏经验而盲目，用无数的邪恶压垮了托付给他们照管的百姓。这正是几乎发生在我身上的事，若非天主迅速地将我从那些危险中解救出来，仁慈地怜惜他的教会和我自己的灵魂。因为，请告诉我，你认为教会中如此巨大的麻烦从何而来？在我看来，我认为唯一的来源是选择与任命主教时的轻率和随意。因为头应当是最强的部分，以便能够调节和控制从身体其他下部升起的邪恶气息；但若它本身软弱，无法抵御那些有害的攻击，它就会变得比实际更软弱，也毁坏了身体的其他部分。为了防止此事发生，天主使我保持在脚的位置上，那是我最初被指派的位置。因为，\Person{巴西略}，除了已经提到的那些，司铎还应具有许多其他品质，但我并不具备；而首要的是这一：他的灵魂应当完全清除对职务的任何贪欲：因为若他对这一尊位有天然的倾向，一旦获得它，就会燃起更强烈的火焰，那人被完全俘虏，将忍受无数的邪恶以稳固地把它保住，甚至到了使用谄媚、或屈从于某些卑鄙可耻之事、或花费大量金钱的地步。因为我现在不谈那些有些人为了争夺尊位而使教会充满谋杀、使城市荒凉的事，免得有人认为我所说的话难以置信。但我的意见是，人们应当如此审慎，以至于逃避职务的重担，一旦担任，若犯了应被罢免的过错，不等他人裁决，而应主动将自己从尊位中开除；因为这样他或许能从天主那里获得怜悯；但若不顾体统地紧握尊位，就是剥夺自己的一切宽恕，甚至点燃天主的怒火，加上第二个比第一个更严重的错误。\par

11. 但没有人能永远承受压力；因为对这尊位的热切渴望，实在是可怕的，真是可怕的。我说这话并不与\Person{保禄}相反，而是完全符合他的话。因为他说什么？\BibleQuote{若有人渴望主教的职务，他就渴望善工。}\BibleRef{弟前 3:1} 我并没有说渴望这\Emph{工作}是可怕的事，而只是渴望权柄和权力。我认为应当用一切可能的热忱将这种渴望从灵魂中驱除，从一开始就不允许被这种情感所占据，以便能够自由地做一切事。因为不渴望被展示拥有这权柄的人，不怕被罢免，并且因不怕就能以基督徒应有的自由做一切事；而那些害怕战栗怕被罢免的人，则忍受着充满各种邪恶的痛苦奴役，常常被迫得罪天主和人。灵魂不应受这种影响；正如在战争中，我们看到那些高尚的士兵愿意战斗并勇敢地倒下，同样，那些获得了这一管理职务的人，应当满足于被祝圣于这尊位或被移除，以基督徒的方式行事，知道这种罢免的报酬不亚于履行职务。因为当任何人为了不做有损于这尊位的事而遭受这类事情时，他为自己招致惩罚那些不义地罢免他的人，并为自己获得更大的报酬。主说：\BibleQuote{几时人为了我的缘故辱骂你们、迫害你们，捏造一切坏话毁谤你们，你们是有福的。你们欢喜踊跃吧！因为你们在天上的赏报是丰厚的。}\BibleRef{玛 5:1}\par

为此，我们必须在各方面保持警醒，并仔细查察，以免此种欲望的火星在某处暗燃。因为我们极愿那些原本不受此情欲支配的人，在担任此职后也能避开它。但若有人在获得尊荣之前，心中已怀藏这可怕而凶残的怪兽，则他获得尊荣后将投入何等烈焰，实难言喻。我本人曾高度怀有此种欲望（请勿以为我会在自我贬抑中说不实之言）：这欲望与其他理由相结合，使我颇感惊惧，促使我逃走。正如迷恋某人者，当被允许接近所爱对象时，因情欲而受更剧烈的折磨；但当他们尽可能远离这些欲望对象时，则驱散了狂热：同样，那些渴望此尊荣的人在接近它时，邪恶变得难以忍受；但当他们不再对此抱有希望时，欲望便随期望一同熄灭。\par

12. 仅此一个动机便非轻小；单凭它本身已足以使我对此尊荣望而却步。然而，如今还需加上另一个不亚于前者的理由。这理由是什么？司铎应当具有清醒的头脑、透彻的辨别力，以及四面八方的无数眼睛，因为他不是为自己一人而活，而是为如此众多的群众。然而，我本人迟钝懈怠，几乎连自己的救恩都难以促成，即使你这因爱我而特别急于隐藏我缺点的人，也会承认这一点。在此请不要向我谈论禁食、守夜、睡硬地及其他刻苦肉身之事：因为你深知我在这些方面多么欠缺；即便我认真操练这些，以我目前的懈怠，它们对于我担任此权威职位也毫无裨益。这些事情对于幽居静室、只关心自己事务的人大有裨益；但当一个人分心于如此众多的群众，并分别进入属下每个人的私虑时，除非他拥有强健而极其刚毅的品格，否则对改善他人能有何显著帮助呢？\par

13. 不要惊讶，在论及此种忍耐时，我还要在灵魂中寻求另一种刚毅的考验。因为对饮食和软榻的漠不关心，我们看到对许多人并非难事，至少对那些生活粗犷、自幼如此教养的人，以及许多其他人而言，身体的锻炼和习惯会使这些劳苦实践的严酷性变得柔和。但侮辱、辱骂、粗言、来自下属的嘲笑（无论是无端还是正当的）、来自统治者和被统治者徒然无谓的责备——这些是很少有人能忍受的，实际上只有极少数人才能承受；人们会看到那些在前述操练中坚强的人，竟被这些事情完全击垮，变得比最凶残的野兽还要狂暴。这样的人我们尤其应排除在司铎职的门槛之外。因为若一位主教不厌恶食物，或赤足行走，对教会的公共利益并无损害；但暴怒的脾气会给他本人及其邻人带来巨大灾难。对于未履行前述事情的人，并无神圣的威胁；但无端发怒的人则受到地狱和地狱之火的威胁。\BibleRef{玛 5:22} 因此，正如热爱虚荣的人一旦掌管众多人，便给火焰添加燃料；同样，那独自或与少数人在一起时无法控制愤怒、轻易被其左右的人，若被托付管理整个群众，便会像被无数折磨者从四面八方刺戳的野兽，自己永不得安宁，且将给托付于他的人带来不可估量的祸害。\par

14. 因为没有什么比不受约束的愤怒（它猛烈急速地冲撞）更能遮蔽理性的纯洁和心智视觉的清晰。有人曾说：\Quote{愤怒甚至摧毁明智的人。} 因为灵魂之眼如同在夜战中变得昏暗，无法分辨朋友与敌人，也不能辨别可敬与卑劣，而是以同样方式依次对待一切；即使必须遭受某些损害，也甘愿忍受一切，以满足灵魂的愉悦。愤怒之火是一种愉悦，它比愉悦更严厉地统辖灵魂，完全扰乱其健全的组织。因为它轻易驱使人走向傲慢、不合时宜的敌意和无理的憎恨，不断使人准备做出无礼而徒然的冒犯；并迫使他们说许多此类之事，灵魂被激情的急流冲走，毫无力量支撑并抵抗如此巨大的冲动。\par

\Person{巴西略}：我再也不能忍受你这样的讽刺了：因为谁不知道你离这种软弱有多么遥远呢？\par

\Person{金口圣若望}：\Quote{那么，我亲爱的朋友，你为何想要把我带到火堆旁，在野兽安静时去激怒它？难道你不知道，我达到这种状态，并非由于任何内在的德行，而是由于我对退隐的热爱？而且，一个如此禀赋的人，当他满足于独处，或只与一两个朋友交往时，他能够逃脱由这种情欲产生的火焰；但如果他陷入所有这些忧虑的深渊，他就不仅会把自己，还会把许多其他人一起拖向毁灭的边缘，并使他们更不关心温和之道。因为一般受管辖的民众往往倾向于将统治者的行为视为一种典范，并使自己效仿他们。那么，当他自己因愤怒而膨胀时，他如何能制止他们的狂怒？当百姓看到统治者易怒时，他们中有谁会立刻渴望变得温和？因为司铎的缺陷是绝对不可能被隐藏的，即使是微小的缺陷也会迅速暴露。就像一个竞技者，只要他待在家里，不与任何人竞争，即使他的弱点非常大，也能掩饰；但当他脱衣上场比赛时，就很容易被识破。同样，对于过着这种隐退和无所作为的生活的人来说，他们的独处成为掩盖他们缺陷的面纱；但一旦被带入公众视野，他们就被迫脱去这种隐退的外衣，并通过他们可见的行为向所有人暴露他们的灵魂。因此，正如他们的正确行为有益于许多人，激发他们同样的热忱，同样，他们的缺点也使人们对实践德行更加漠不关心，并鼓励他们在追求卓越时懈怠。因此，司铎的灵魂理应在每一方面都闪耀美丽，以便能够使那些看见它的人的灵魂喜悦和启迪。因为普通人的过失，几乎是在黑暗中犯下的，只毁掉那些行恶的人；而一个处于显赫地位、为众人所知的人的错误，则对所有的人造成共同的伤害，使那些已经跌倒的人在向善的努力中更加懈怠，并使那些想要谨慎自守的人陷入绝望。此外，无关紧要的人的过失，即使暴露，也不会对任何人造成值得一提的伤害；而那些占据最高荣誉席位的人，首先在所有人面前显而易见，如果他们犯了最小的错，这些小事在他人看来也是巨大的；因为所有的人都衡量罪恶，不是根据冒犯的大小，而是根据犯人的地位。因此，司铎理应以一种金刚石般的盔甲从四面八方护卫自己，以极大的认真和对他生活方式的持续警惕，免得有人发现一个暴露和疏忽的地方，并给予致命的一击：因为所有环绕他的人都准备打击和推翻他；不仅敌人和反对者，甚至许多自称朋友的人也是如此。}\par

因此，被选为司铎的人的灵魂理当被赋予天主恩宠赐予那些被投入巴比伦火窑的圣人们身体的那种力量。\EditorialNote{达尼尔书 3} 柴捆、沥青和麻屑不是这火的燃料，而是更可怕的东西：因为他们所遭受的不是物质之火，而是嫉妒的吞噬一切的火焰从四面升起，袭击他们，对他们的生命的考验比当时那火对青年身体的考验更彻底。因此，当它发现一点茅草的痕迹时，就迅速附着上去；并且这有缺陷的部分完全被烧尽，而其余的部分，即使比阳光更灿烂，也被烟雾烧焦和熏黑。因为只要司铎的生活在各方面都良好地规范，他就不会受到阴谋的伤害；但若他不幸忽略了一些小事（这是人之常情，因为他在人生险恶的海洋中航行），那么他其他所有的善行都无助于他逃脱控告者的口舌；那个小错误掩盖了一切其他优点。并且所有人都准备对司铎作出审判，仿佛他不是血肉之躯，或不是承继人性，而是像天使一样，免除了各种软弱。正如所有人当暴君强大时都害怕和奉承他，因为他们无法推翻他；但一旦看到他的事务不顺利，那些不久前的朋友就放弃虚伪的尊敬，突然变成他的敌人和对手，发现了他所有的弱点，就攻击他，并废黜他的统治；同样，对于司铎也是如此。那些不久前在他强大时尊敬和讨好他的人，一旦找到一点把柄，就准备废黜他，不仅像对待暴君，而且比那更可怕。正如暴君害怕他的卫兵，同样司铎最惧怕的是他的邻人和同工。因为没有别人像他们这样觊觎他的尊位，或如此了解他的事情；一旦有什么事发生，由于近在咫尺，他们比别人更先察觉，即使他们诽谤他，也容易让人相信，并且通过夸大小事，将猎物捕获。因为宗徒的话被颠倒了：\BibleQuote{无论一个肢体受苦，所有肢体都一同受苦；或一个肢体受尊荣，所有肢体都一同欢乐。}\BibleRef{格前 12:26} 除非一个人能以极大的谨慎抵挡一切。\par

那么，你当真要派遣我投入如此巨大的战斗吗？你以为我的灵魂能胜任这般形态各异、变化多端的较量吗？你从哪里学到这个，又是从谁那里听来的？如果天主向你证实了此事，请将神谕指示给我，我便服从；但若你不能，而仅凭人的意见作判断，那么请你摆脱这种错觉。因为关于我的事，信任我比信任别人更为公平；诚如经上说：\BibleQuote{除了人内里的神，没有人知道人的事。}\BibleRef{格前 2:11} 倘若我接受了这职务，便会使自己和选举我的人成为笑柄，并要遭受巨大损失才回到我现在的这种生活状态。我相信，即使以前未能说服你，现在通过这番话也已使你信服了。因为不仅是恶意，还有一种更为强烈的东西——对这尊位的贪欲——常会驱使许多人武装起来反对占有它的人。正如贪财的儿女受父母年老的压制，同样，这些人见到任何人长期持有司铎职务时——既然消灭他是一种恶行——便急忙将他废黜，因为个个都渴望取而代之，且都期待那尊位会转移到自己身上。\par

15. 你愿意我向你展示这场斗争的另一个充满无数危险的方面吗？那么，请来看一看通常在选举教会尊位时的公共节庆，你便会看到司铎受到与他所治理的人数一样多的指控。因为所有享有授予荣誉特权的人那时都分裂成许多派别；你永远找不到长老会议意见一致，或关于赢得尊位的人意见一致；而是彼此分离，一人偏爱这个，另一人偏爱那个。原因在于他们并非都注目于同一件事——那应是唯一注意的对象，即品格的卓越；而是举出别的资格来推荐此荣誉；例如，有人说：\Quote{让他当选，因为他出身显赫家族}，另一人说：\Quote{因为他拥有巨大财富，无需靠教会的收入来维持}，第三个人说：\Quote{因为他从敌营那边过来}；有人急于将优先权给予与他亲密的人，另一人给予与他有血缘关系的人，第三人给予谄媚者，却无人关注真正合格的人，或对其品格作任何考验。我远非认为这些是判定一个人是否适合司铎职务的可靠标准，以至于即使有人显出极大的虔诚——这在此职务的履行中并非小帮助——我也不敢仅凭此就认可他，除非他碰巧将良好的才能与虔诚结合在一起。因为我认识许多人，他们对自己实行了恒常的克制，并以斋戒消耗自己；只要他们被允许独处，只关注自己的事，他们便蒙天主悦纳，且日复一日在此类学识上有了不小的长进；但一旦他们进入公共生活，被迫纠正群众的愚昧时，其中有些人从一开始就证明无法胜任如此重大的任务，另一些人在被迫坚持时，则放弃了从前严格的生活方式，从而给自己造成了巨大伤害，却对他人毫无裨益。如果有人将全部时间花在最底层的圣职上，并活到极高的老年，我不会仅出于对他年岁的尊敬就将他提升到更高的尊位；因为假如到了那个年纪，他仍然不适合此职务，那又如何？我此刻这样说，并非有意羞辱白发，也不是立下一条律法绝对将来自隐修圈子的人排除在此权威之外（因为有许多例子表明，从那里出来的人，在此尊位上显得光辉灿烂）；但我急于证明的是：如果虔诚本身加上年高都不足以表明一个获得司铎职务的人确实配得上它，那么以前所举的理由也几乎无法做到这一点。还有另一些人提出更加荒谬的借口；有些人是为了不让自己成为反对者的一员而被列入圣职人员的行列，另一些人则是因为他们的恶性，生怕若被忽视会做大的坏事。还有什么比这更违背正规呢？即恶贯满盈的人竟因那些本应受罚的事而受到奉承，并因那些本应被禁止进入教会门槛的事而升到司铎的尊位。那么，当我们把如此神圣可畏的事物交托给邪恶或无用的人去玷污时，还需要寻找天主发怒的原因吗？因为当一些人被托付管理完全不适合他们的事物，而另一些人被托付管理超出他们天生能力的事物时，他们使教会变得像\OriginalTerm{尤里普斯海峡}{Euripus}一样。\par

从前我常讥笑世俗的统治者，因为他们在分赐荣誉时，不是以道德卓越为准则，而是以财富、年资和人的等级为准则；但当我听说这种愚妄也侵入了我们的事务时，我便不再认为他们的行为如此可憎了。这有什么稀奇呢？世俗之人喜爱群众的赞美，为利益而行事，他们犯这些罪，而那些至少装作不受这些影响的人，却一点也不比他们更好，虽然从事于属天事物的竞赛，行为却仿佛所提交裁决的问题关乎几亩田地或类似的事。因为他们随手选取平庸之人，委派他们管理那些事——为了这些事，天主的独生子并未拒绝空虚自己，舍弃祂的荣耀，降生成人，取了奴仆的形体，受人唾污，受人掌击，在肉身中受辱而死。他们甚至不止于此，还在这罪上加添更严重的罪行；因为他们不仅选举不配的人，而且还驱逐合格的人。仿佛必须从两面毁坏教会安全，又似乎先前的挑衅不足以激起天主义怒，他们又策划了另一件同样有害的事。因为我以为驱逐有用之人与强行塞入无用之人同样可憎。这事实在发生，以致基督的羊群在任何方向都找不到安慰，也无法自由呼吸。难道这些行为不该受一万次雷霆之击，以及比圣经所警告的更炽热的地狱之火惩罚吗？然而，这些极恶之事却被祂容忍着，祂不愿罪人死亡，而愿他悔改而存活。谁能充分惊叹祂的慈爱，惊讶于祂的慈悲呢？光荣归于祢，主！光荣归于祢！祢慈爱的深渊何其浩瀚！祢宽容的富饶何其伟大！那些因祢的名从卑微和隐晦升到尊荣显位的人，竟用所得的尊荣反对赐予者，行狂妄僭越之事，侮辱圣物，拒绝并驱逐热忱之人，好让恶人极其安全、毫无畏惧地随意毁坏一切。如果你想知道这可怕之恶的原因，你会发现它们与先前所述相似；因为它们有一个共同的根和母亲，可以说——就是嫉妒；但这嫉妒以几种不同形式表现出来。因为我们听说，一个是因为年轻而被从候选人名单中剔除；另一个是因为不懂奉承；第三个是因为得罪了某某人；第四个是因为怕某某人看见他所举荐的人被拒绝而此人当选会感到痛苦；第五个是因为仁厚温和；第六个是因为令罪人畏惧；第七个出于其他类似原因；因为他们毫不费力地找到所需的各种借口，甚至在没有其他借口时，以某人财富丰裕作为反对理由。事实上，他们能够发现其他理由，要多少有多少，证明一个人不应突然被提拔到此尊荣，而应温和渐进。在此，我想问一个问题：\Quote{那么，那位必须与如此风暴搏斗的主教该如何行？他如何能抵挡如此巨浪？他如何能击退所有这些攻击？}\par

因为如果他秉公处理，那些与他及候选人敌对和反对的人，便专事争斗，日日挑起争端，对候选人堆砌无尽的轻蔑，直到他们使候选人被除名，或引入自己的宠信。事实上，这就像一位舵手，有海盗与他同船，时时每刻都在密谋对付他、水手和士兵。如果他宁愿讨好这些人而不顾自己的救恩，接受不配的候选人，那么他就会有天主作他的敌人来代替他们；还有比这更可怕的吗？但他与他们的关系将比以前更加尴尬，因为他们会彼此联合，从而变得比以前更强大。因为正如猛烈的风从相反方向碰撞时，原本平静的海洋突然狂暴起来，掀起浪头，毁灭航行其上的人，同样，当教会接纳了腐坏之人，它原本平静的表面便覆盖着汹涌的浪涛，布满船难。\par

16. 那么，请想一想，一个人应当是怎样的人，才能抵抗如此的风暴，并巧妙地应对这些危害公共福祉的巨大障碍；因为他应当庄严而不傲慢，可畏而和蔼，善于发令而平易近人，公正而有礼，谦卑而不卑屈，刚强而温柔，这样他才能成功克服所有这些困难。他应当以大权威举荐真正适合此职务的人，即使所有人都反对；同样以大权威拒绝不适合的人，即使所有人都合谋支持他；他应只注视一个目标，即建立教会，不受任何敌意或偏袒的驱动。那么，你现在认为我拒绝这项职务的职责是合理的吗？但我尚未向你道尽所有理由；因为我还有一些别的要提及。请不要厌倦听一位友好而真诚的人说话，他希望能为自己辩护，使你免于责备；因为这些话不仅对你要为我所作的辩护有益，而且对妥善履行这项职务也不无小补。因为即将踏上这条路的人，在着手从事这项职务之前，必须彻底探究一切。你问为什么？因为清楚知道一切的人，至少会有这样的好处：当这些事发生时，他不会感到陌生。那么，你是否愿意我先处理关于照顾寡妇的问题，或是关于照料童贞女的问题，或是关于司法职务的困难？因为在这些情况中，每一种都有不同的焦虑，而恐惧更甚于焦虑。\par

首先，从看似比其他议题更简单的问题开始：照顾寡妇一事似乎仅在花费金钱上使照管者忧虑；但实情并非如此。当她们需要被登记入册时，仍需审慎查考，因为不加甄别地将她们列入名册已造成无穷祸害。她们曾毁坏家庭、拆散婚姻，且常被发现在偷窃、侵占及诸如此类不体面的事中。若用教会的收入供养这样的女子，必招致天主的惩罚、在人前遭受极严厉的谴责，并使那些乐意行善的人的热忱减退。因为谁愿意将他奉基督之名所应施舍的财富，浪费在那些诋毁基督之名的人身上呢？因此，必须进行严格而准确的甄别，以免穷人的供给予被上述妇人以及有能力自养的人浪费。在这番甄别之后，还有另一类并非微小的忧虑：要确保供应如同从泉源涌出一样丰盈不断；因为贫穷是一种贪得无厌的恶，好抱怨且忘恩负义。需要极大的谨慎与热忱，才能堵住抱怨者的口，剥夺他们的任何借口。多数人见到某人超脱于贪财，便立刻认为他适合担任这管家的职务。但我以为，单有这种慷慨之心永远不够，尽管它应优先于一切其他品质；若没有它，人将成为毁灭者而非保护者，是豺狼而非牧者；然而，除了这一品质，还须要求具备另一种品质。这品质便是容忍，这是人心中一切美善之因，它仿佛推动并引领灵魂进入宁静的港湾。因为寡妇这类人，因其贫穷、年迈及本性，常放纵无度的言语（我最好这样称呼）；她们不合时宜地喧嚷、无谓地抱怨哀叹那些本应感恩之事，并对那些本应欣然接受的事提出指控。监管者应以宽厚的精神忍受这一切，不为她们过分的烦扰或无理的抱怨所激怒。因为这类人值得因不幸而被怜悯，不应受侮辱；践踏她们的苦难，在贫困带来的痛苦上再加侮辱的痛苦，是极端的残忍。因此，一位最智慧的人，洞察人性的贪婪与骄傲，又思量贫穷的本性及其可怕的力量——使最高贵的品格也消沉，并常诱使人做出同样无耻的行为——为使人不被控告所激怒，也不因持续的纠缠而成为他本应援助者的敌人，便教训他要和蔼可亲、易于接近求告者，说：\BibleQuote{你当向穷人侧耳，以温和的言语回答他。}\BibleRef{德 4:8} 他略过那足以令人恼怒的事（因为对那被克服的人还有什么可说呢？）而是对那能忍受他人软弱的人说话，劝勉他在施予之前，先以温和的面容与柔和的言语纠正求告者。但若有人虽不侵占（寡妇的）财产，却以无数斥责、侮辱及恼怒对待她们，他的施舍非但不能减轻贫穷所生的沮丧，反而以辱骂加重了痛苦。因为她们虽可能因饥饿之迫而不得不极其无耻，但她们对此强迫也感到痛苦。因此，当她们因饥饿的恐惧而被迫乞讨，因乞讨而被迫放弃羞耻，又因无耻而受辱骂，沮丧之力便变得复杂，且带着深重的阴郁沉积在灵魂上。负责照顾她们的人应当如此忍耐，不仅不以自己的怒气增加她们的沮丧，还要以其劝勉除去其中大部分。因为正如受辱者虽享有丰裕，但因所受侮辱之打击也感不到财富的益处；同样，得温和言语相待，且礼物伴随鼓励之人，便越发欢欣雀跃，所赠之物因赠予的方式而倍增其价值。我说这话并非出于自己，而是借用了刚才引用其箴言的那一位：\BibleQuote{我儿，不要破坏你的善行，施舍时不要说令人不安的话。露水岂不消散炎热？话语胜过礼物。看啊！话语岂不胜过礼物？但两者都与有恩惠的人同在。}\BibleRef{德 18:15}\par

监管这些人的人不仅应当温和容忍，还应善于管理财产；若缺少此资质，穷人的事务又将陷入同样的困境。不久前有一位受托此职务的人，积聚了大量钱财，既不自己使用，也极少用于需要的人，而是将大部分埋藏在地下，直至灾祸临头时，全部落入敌人之手。因此需要极大的远见，使教会的资源既不丰裕有余，也不匮乏不足，而是一切所备的供养迅速分给需要的人，并将教会的宝藏储存在她治下之人的心中。\par

此外，在接待旅客与照顾病患的事上，请想一想需要多大的花费，以及主管这些事务的人需要何等精确与辨别的能力。因为这项花费往往必须比刚才我所说的更大，主管者应当将谨慎与智慧与供给之道的技巧结合起来，以便使富裕者竞相慷慨施舍，这样在提供病患所需时，不会困扰那些施舍者的灵魂。但在此更需表现出更高的热心与勤勉，因为病患是难以取悦的人，容易倦怠；若不运用极大的精确与照管，即使一个微小的疏忽也足以给病人造成巨大的损害。\par

17. 但在看守贞女的事上，恐惧更大，因为所拥有的更宝贵，这群羊也比其他的更高贵。确实，已有无数满盈无数恶习的女子冲入了这群圣者的行列；在此，我们的忧伤比在其他事上更大；因为一个失足的贞女和一个失足的寡妇之间的差别，正如一个自由出身的少女和她的婢女之间的差别。对于寡妇，轻浮、彼此辱骂、谄媚或厚颜、到处在公众前露面、在集市上闲逛，已成为常事。然而贞女追求更高尚的目标，热心寻求最高超的哲学，并宣称要在世上显现天使所过的生活，在肉身之中立志去做属于无形体能力的事。此外，她不应作许多或不必要的旅行，也不可说闲话和随意的话；至于辱骂和谄媚，她甚至不应知道这些名字。为此，她需要最谨慎的守护和更大的帮助。因为圣德的仇敌总是出其不意地埋伏等候这些人，随时准备吞吃她们中任何一个如果滑倒失足；也有许多男人为她们设下网罗；除这些之外，还有她们人性本身的激情，所以一般而言，贞女必须为一场两面的战争装备自己：一面从外攻击她，另一面从内压迫她。因这些缘故，那监管贞女的人深陷恐慌，而若有任何违背他意愿的事发生——愿天主禁止——危险和痛苦就更大。因为\BibleQuote{一个幽居的女儿，能使父亲失眠，因怕她无子，或虚度青春（未嫁），或被（丈夫）憎恨，而忧虑使他不能入睡}，\BibleRef{德 42:9} 那么，那人的忧虑不关乎这些事，而关乎远为更大的事，他又将受何等的苦呢？因为在此，被拒绝的不是一个男人，而是基督自己；这种不生育也不仅是可耻的，而是恶事以灵魂的丧亡告终；因为经上说：\BibleQuote{凡不结好果子的树，就砍下来，丢在火里。}\BibleRef{玛 3:10} 而被神圣新郎休弃的人，仅领受休书离去是不够的，还要受永罚的惩罚。此外，按肉身之父有许多事使他对女儿的看守变得容易；因为母亲、乳母和许多婢女共同帮助父亲保全女儿。因为她既不被允许经常匆忙进入集市，即便去了，也不必向任何一个行人显露自己，夜色如同屋宇的墙垣一样，能遮蔽不愿被看见的人。除这些外，她也免除了所有可能迫使她与男人目光接触的缘由；因为对生活必需品的忧虑、压迫者的威胁，或任何这类事，都不会使她陷入这种不幸的必要性，她的父亲在这些事上代她处理；而她本人只有一件忧虑，就是避免做或说任何有损于她应有的端庄品行的事。但在另一种情况（指看守贞女）下，有许多事使看守贞女变得困难，甚至对父亲而言是不可能的；因为他不能让她与自己同住一屋，因为那样同居既不体面也不安全。即使他们自己没有受损，仍保持纯洁无瑕，他们也要为那些因他们而跌倒的灵魂交账，如同他们彼此犯罪一样。由于不可能同住，便难以了解性格的动向，也难以抑制失序的冲动，或训练和改善那些更有秩序、更协调的冲动。干预她外出的习惯也不是容易的事；因为她的贫穷和缺乏监护人，使他无法精确查究她行为是否得体。因为她不得不处理自己的一切事务，有许多外出的借口，至少如果她不愿自律的话。那命令她总待在家的人，理应切断这些借口，为她提供生活所需的独立，并给她一个管理这些事务的妇女。他还必须使她远离葬礼和夜间的节庆；因为那条狡猾的蛇太懂得如何借着善行散布它的毒液。贞女必须受到四面防护，全年极少出门，除非被无法抗拒的必要所迫。若有人说这些事没有一件是主教应当操持的，让他明白，关乎每件发生之事的一切忧虑和缘由，都必须向他汇报。而且，由他亲自管理一切，从而免于因他人的过失而不得不承受的抱怨，远比他不管理，却因他人所行之事而恐惧被追究更为有利。此外，由他亲自办理这些事，他能非常轻松地完成全部；但若他被迫透过说服每个人的意见来办理，那么他因那些反对他、与他的决定争斗的人所遭受的烦恼和纷扰，远比他免于单枪匹马劳动所得的喘息为大。然而，我无法一一列举与看守贞女有关的一切忧虑；因为当她们要登记录名时，会给受托于此的人带来不小的麻烦。\par

再者，主教职务的司法部门涉及无数烦恼、大量时间消耗，以及远超审理世俗事务之人所经历的困难；因为发现准确的公义是件苦差事，即使找到了，也很难避免毁掉它。而且不仅遭受时间和困难损失，还有不小的危险。因为至今，有些较软弱的弟兄因涉入事务，因未得到庇护，在信德上遭遇了船难。因为许多遭受不公的人，与施不公的人一样，仇恨那些不帮助他们的人，他们不会考虑所涉事项的复杂性、时势的困难、司铎权柄的限度或任何此类事情；他们是无情的审判者，只承认一种辩护——从压迫他们的邪恶中解脱出来。不能提供这种解脱的人，即使提出无数借口，也永远无法逃脱他们的谴责。\par

谈到庇护，让我揭示另一个吹毛求疵的借口。因为如果主教不像城里的闲人那样每天都进行一轮拜访，就会引起难以言喻的冒犯。因为不仅是病人，连健康的人也渴望被探望，这并不是出于虔诚促使他们如此，而是在大多数情况下，他们假装要求尊敬和殊荣。如果他曾经因为某种需要，为了教会的共同利益，更频繁地拜访某位较富有、较有势力的人，他立刻就被贴上阿谀奉承的标签。但我为何只谈庇护和拜访呢？仅仅从他们与人的打招呼方式，主教们就得忍受如此多的指责，以至于常常因沮丧而压抑和淹没；事实上，他们还要接受对其眼神使用的审视。因为公众严厉批评他们最简单的动作，注意他们声音的语调、面容的表情和笑容的程度。有人会说：他对那人开怀大笑，满面春风地打招呼，声音清晰，而对我只说了简短敷衍的话。在大型集会中，如果他交谈时没有将眼睛转向每个方向，大多数人就宣称他的行为是侮辱性的。\par

那么，除非他极其坚强，谁能应对如此多的控告者，要么完全避免被起诉，要么如果被起诉就能逃脱？因为他必须要么没有任何控告者，要么如果这不可能，就洗清所提出的指控；如果这又不那么容易，因为有些人喜欢提出虚妄无理的指控，他必须勇敢地对抗由这些抱怨产生的沮丧。确实，被公正指控的人可能容易容忍控告者，因为没有比良心更苦毒的控告者；因此，如果我们先被这个最可怕的对手抓住，我们就能更容易忍受外在的较温和的对手。但那些良心无愧的人，当受到空洞的指控时，会迅速激起愤怒，并容易陷入沮丧，除非他事先练习过如何忍受大众的愚昧。因为对于一个毫无理由地被诬告和定罪的人来说，面对如此大的不公，避免感到一些烦恼和恼怒是完全不可能的。\par

什么时候需要将某人从教会的完整共融中切断，主教所受的痛苦又怎能尽述？但愿这祸患止于痛苦！但实际上，损害并不小。因为恐怕那人若受了超过应得的惩罚，就会经历蒙福的保禄所说的话，并\Quote{被过度的忧愁所吞没。} \BibleRef{格后 2:7} 因此，在这件事上也需要最精确的准确性，以免本应有益的事反而成为他更大损害的缘由。因为无论他在这种处理方式之后犯什么罪，每项罪所激起的义怒都必须由那位如此不熟练地将刀切入伤口的医生来分担。那么，一个人不仅要为自己单独犯的过犯交账，还要因别人所犯的罪陷入极度的危险，他将遭受何等严厉的惩罚？因为如果我们因自己的恶行颤栗受审，相信自己无法逃脱来世的火，那么必须为那么多别人负责的人，又该期待承受什么呢？为证明此事的真实性，请听蒙福的保禄——更确切地说，不是听他，而是听在他内说话的基督的话：\Quote{你们要服从那些管理你们的人，并要顺服，因为他们为你们的灵魂警醒，好像那些将交账的人。} \BibleRef{希 13:17} 这威胁的恐惧能小吗？不可言说：但这些考量足以使最不信和顽固的人相信，我如此逃脱并非出于骄傲或虚荣，而仅仅出于对自身安全的恐惧和对职务严重性的考虑。\par

\SourceRecord{Translated by W.R.W. Stephens.；Nicene and Post-Nicene Fathers, First Series；Vol. 9.；Edited by Philip Schaff.；Buffalo, NY: Christian Literature Publishing Co.,；1889.；<http://www.newadvent.org/fathers/19223.htm>.}

% structure: p0001 source=h1 target=section reason=page-title

\section{\WorkTitle{论司铎职（卷四）}}

\Person{巴西略}听了这话，稍作停顿，便这样答道：\par

如果你自己渴求这个\OriginalTerm{职分}{office}，你的惧怕本是合理的；因为人若渴求承担此职，便是承认自己有能力管理它，而一旦受托付后失败，他便不能以缺乏经验为托词，因为他早已事先剥夺了这一借口——仓促地抓住了\OriginalTerm{职务}{ministry}；任何自愿而有意接受此职的人，再也不能说：\Quote{我在这方面违背我的意愿犯了罪——违背我的意愿我毁坏了某个灵魂；}因为那将来审判他的，要对他说：\Quote{既然你自知如此缺乏经验，没有能力承担此事而不受责难，为何你如此急切、如此鲁莽，竟去经办远超你能力的事？谁强迫你这样做？你曾退缩或逃避吗？难道有人强行拉你来吗？}但你将听不到这样的话，因为你没有什么可自责的；所有人都清楚你丝毫不渴求这一尊荣，因为事成全赖他人所为。因此，那些渴求此职分的人犯错时毫无宽赦的机会，而你的处境则为你提供了充分的借口。\par

\Person{金口若望}听到这话，我摇了摇头，微微一笑，钦佩此人的天真，便这样对他说：我本希望事情如你所说，最卓越的人，但这并非为了使我能接受我最近逃避的职分。因为，即使我毫不思索、毫无经验地接受照料基督的羊群不会受到惩罚，但对我来说，受托如此重大的责任后，向那托付我的人显得如此卑劣，比任何惩罚更痛苦。那么，我为什么希望你没有看错呢？实在是为了那些可怜而不幸的人（因为我必须这样称呼他们，他们不知道如何善尽此职的职责，尽管你说一万次他们是迫不得已接受此职，因此他们的错误是出于无知的罪）——为了这些人，我再说，使他们得以逃脱那\BibleRef{玛 25:30}不灭的火、外面的黑暗和\BibleRef{谷 9:44}不死的虫，以及被腰斩的刑罚，并与假善人一同丧亡。\par

但我能为你做什么呢？事情并非像你所说，绝不如此。如果你愿意，我可以从国度的例子中给你一个我所坚持的证明——国度在天主眼中不如司祭职重要。\Person{克士}的儿子\Person{撒乌耳}本人毫不渴求成为君王，他正在寻找母驴，来向先知询问它们。然而，先知却向他谈及王权，但即使如此，他也没有贪婪地追逐它，尽管从先知那里听说，却退缩并推辞道：\Quote{我是谁，我的父家是什么？}\BibleRef{撒上 9:21}那么，后来他滥用天主赐予他的尊荣时，那些话能否救他脱离立他为君者的愤怒？当他受到\Person{撒慕尔}责备时，他能说：\Quote{我难道贪婪地追逐王权和主权吗？我本想过平凡人宁静和平的生活，但\Emph{你}却把我拖到这个显赫的职位。若我仍处于卑微的身份，我本可轻易避开所有这些绊脚石，因为若我是无名大众之一，我决不会被派去远征，天主也不会将攻击\Place{阿玛肋克人}的战争交在我手中，若没有交在我手中，我就不至于犯这罪了。}但所有这些辩解都软弱无力，不仅软弱，而且危险，因为它们反而激起天主的义怒。因为那被天主提拔到高位的人，绝不能以他的尊荣伟大作为他错误的借口，而应将天主的特别恩宠作为进一步进步的动机；而那些因为获得了某种非凡的尊荣便以为自己可以自由犯罪的人，岂不是在表明天主的慈爱成为了他个人过犯的原因？这总是那些过着无神而粗心大意生活之人的论调。但\Emph{我们}绝不可如此思想，也不可堕入这些人的疯狂愚昧，而要时常努力善用我们所有的能力，在言语和思想上保持虔诚。\par

因为（撇开王国不谈，回到我们谈论的更直接的主题——司祭职：）厄里并不贪求获得他的高位，然而当他在这职位上犯罪时，这对他有何益处呢？但我为何说\Quote{获得}呢？即使他愿意，他也无法避免，因为他有法律义务接受此职。他属于肋未支派，必须承担从他祖先传下来的那个高位，尽管如此，他仍因他儿子的不法行为而付出了不小的代价。至于犹太人的第一位大司祭，天主曾对梅瑟说了许多话关于他，当他独自无法抵挡如此众多群众的狂热时，若不是他兄弟的代祷平息了天主的义怒，他岂不几乎被消灭？\BibleRef{出 32:10} 既然我们提到了梅瑟，最好也能从他身上发生的事来证明我们所说的是真理。因为这位圣梅瑟如此远离贪求犹太人的领导，以至于他拒绝了这个提议（\BibleRef{出 4:13}），并在天主命令他接受时推辞，从而激怒了任命他的天主；不仅那时，而且后来当他开始治理时，他宁愿死去以求摆脱这职责，他说：\Quote{杀了我吧，} 他说，\Quote{如果你要这样待我。} 结果如何呢？当他在默黎巴水边犯罪时（\BibleRef{户 20:12}），这些屡次的推辞能为他开脱吗？它们能说服天主饶恕他吗？他为何被剥夺了预许之地？没有别的原因，正如我们所知，就是因他的这个罪，使那位非凡的人被禁止享受他所治理之人所获得的同样福分；但在许多劳苦和磨难之后，在那难以言表的漂泊之后，在如此多的争战和胜利之后，他死在了那地之外，为了到达那地他曾经历如此多的劳苦和试炼；尽管他经历了深海的暴风骤雨，最终却未能享受港湾的福乐。由此你便看出，不仅那些贪求此职的人对所犯的罪行无可推诿，就连那些因他人的野心而担任此职的人也是如此；因为实在，那些由天主亲自拣选担任此高位的人，即使他们屡次拒绝，仍受到如此严厉的惩罚，没有任何事物能拯救他们中的任何一人脱离这危险——无论是亚郎，厄里，还是那位圣人、先知、行奇迹者、在地上众人中最谦和的（\BibleRef{户 12:3}）、与天主面对面说话如同人与朋友交谈（\BibleRef{出 33:11}）的伟大人物——我们这些远远不及那伟人德行的人，几乎不能以我们从未贪求过这一尊荣的自觉作为充分的借口，尤其当今许多的任命并非出于天主的恩宠，而是出于人的野心。天主拣选了犹达斯，将他算作神圣团体中的一员，并将宗徒职务的尊荣托付给他，如同其他人一样；是的，他还给了他一些超过他人的东西——管理钱财的责任（\BibleRef{若 12:6}）。但这又怎样呢？当他后来滥用这两种信托，背叛了他奉命宣讲的那位，并挪用了他应善为管理的钱财时，他逃脱了惩罚吗？不，正因如此，他反而给自己招致了更大的惩罚，这很合理。因为我们不可利用天主赐予我们的崇高尊荣来冒犯祂，而应以此更令祂喜悦。但一个人因为自己比别人得到更高的尊荣，就声称自己在应受罚时可以豁免，这与那些不信的犹太人类似，他们听到基督说：\Quote{我若没有来，没有对他们说话，他们就没有罪；} \Quote{我若没有在他们中间行过别人未曾行的事，他们就没有罪}（\BibleRef{若 15:22}），竟反过来责备人类的救主和恩主说：\Quote{那么，你为什么来？你为什么说话？为什么行奇迹？是为了更加重我们的惩罚吗？} 但这乃是疯狂和完全无知的话。因为大医师来不是要放弃你，而是要医治你——不是在你生病时忽略你，而是要彻底治愈你的疾病。但你却自愿从祂手中退出；因此你要接受更严厉的惩罚。因为正如你若听从祂的治疗就能从先前的疾病中得释放，同样，如果你看到祂来帮助你却逃离了祂，你将再也不能洗净这些软弱，而既然你不能，你就要既为这些疾病受罚，又因为你使天主对你的眷顾归于无效而受罚。因此，我们这样行事的人，在蒙天主赐予尊荣之前与之后所遭受的折磨并不相同，而是之后比之前远为严厉。因为那甚至受了善待仍不成为善的人，当受更严厉的惩罚。既然你的这个借口已经被证明是站不住脚的，不仅不能拯救那些以此为庇护的人，反而使他们更加暴露，我们就必须为自己提供其他的安全途径。\par

\Person{巴西略}：请告诉我，那是什么样的安全途径？因为就我而言，你刚才的话已使我几乎不能自制，陷入如此的恐惧和战栗之中。\par

\Person{金口圣若望}：我恳求你，不要这样沮丧。因为对于我们这些软弱的人，安全在于根本不担任此职；而对于你们这些强壮的人，安全在于将救恩的希望——除天主的恩宠外——寄托于避免任何有负于这恩赐及赐予它的天主的行为。因为那些凭自己的野心获得此尊位后，或因懒惰、或因邪恶、甚至因缺乏经验而滥用职分的人，确实应受最严厉的惩罚。但这并非暗示那些没有如此野心的人就有赦免。依我之见，即便有千万人恳求敦促，人也应先省察自己的内心，仔细考量整件事，然后才屈从于他们的强求。没有人敢在不是建筑师时承担建造房屋的事，也没有人会在不是熟练医生时尝试医治有病的身体；即使多人催促，他也会推辞，并不耻于承认自己的无知。那么，将要受托照管如此众多灵魂的人，难道不该事先省察自己吗？他会因这个命令，那个强迫，或害怕得罪第三人，就接受这职务——即使他是最无经验的人？若是如此，他怎能不将自己和他们一起推入明显的痛苦？如果他保持原状，或许还能得救；但现在，他却使他人与自己一同灭亡。因为他从何希望救恩？从何获得宽恕？那时谁能成功为我们代祷？也许是现在催促并强行拉我们的人？但那时谁能拯救这些人？因为连他们也需要代祷，才能逃脱烈火。我说这些不是要吓唬你，而是如实陈述事情真相。请听圣宗徒\Person{保禄}对他亲爱的门徒、他的爱子\Person{弟茂德}所说的话：\BibleQuote{不可轻易给人覆手；不可在别人的罪上有分。}\BibleRef{弟前5:22}你难道没有看见，我们——就我们能力所及——使那些准备推举我们担任此职的人免于多大的责备乃至惩罚吗？\par

2. 因为对于那些被选的人，仅仅以\Quote{我没有自邀担任此职，也无法避免我未曾预见的后果}为借口是不够的；同样，那些祝圣他们的人也不能以不知道被祝圣者为由作为充分的辩护。对他们而言，因对被推举者无知而受的指控反而更大，而看似为他们开脱的，反倒更成了控告。因为这是多么荒谬的事！那些想买奴隶的人，会带他去见医生，要求卖方提供担保，并向邻居打听消息，经过这一切，还不立即冒险购买，而是要一段试用期；而那些将要接纳任何人担任如此重大职务的人，却轻率、漫不经心地给予证明和许可，不加进一步查究，仅仅因为某人的愿望，或为了讨好某人，或为了避免得罪某人。在那一天，当那些本该为我们辩护的人自身都需要辩护者时，谁能成功为我们代祷？因此，将要祝圣的人应当勤加查问，而将要被祝圣的人更应如此。因为虽然祝圣他的人要分担他在职期间所犯任何罪过的惩罚，但他本人却非但不能逃脱惩罚，反而要付出比他们更重的代价——除非选他的人出于某种属世的动机，违背了他们自己认为合理之事。如果他们被发现这样做，明知一个人不相称，却因某种借口推举他，那么他们所受惩罚的分量将与他相等，甚至指派不相称之人的人所受惩罚更重。因为授权给一个存心毁灭教会的人，肯定要为此人犯下的暴行负责。但如果他没有这样的过犯，只是被别人意见误导，即使如此，他也绝不能完全免罚，只是惩罚略轻于被祝圣者。那么，怎么办呢？选举者可能被虚假报告所骗。但被选者不能说\Quote{我对自己无知}，如同他人对他无知一样。因此，既然他将承受比推举他之人更严厉的惩罚，他就应当比他们对自己作更仔细的省察。如果他们因无知而强拉他，他就应当自己出面，就任何能使他们免于受骗的事细心教导他们；这样，在表明自己不值一试后，他就能逃避如此崇高职务的重担。\par

因为在战争、商业、农业及今生其他部门的技艺中，当提出某项计划时，农夫不会承担驾船之责，士兵不会耕种田地，领航员不会率领军队，否则便有万死之险，原因何在？岂不正是这样：每个人都预见到自己无能将导致的危险？那么，当损失涉及琐事时，我们竟如此深思熟虑，拒绝向强迫的压力让步；而当惩罚是永恒的，就像对于那些不知道如何处理司铎职的人一样，我们却只愿轻率地投身于如此巨大的危险，然后以他人的迫切恳求作为借口？但那位将来审判我们的主，不会接受这样的辩词。因为我们在超性之事上当远比在肉性之事上更为谨慎。但现在我们却未见表现出同等的谨慎。请告诉我：假设有一个人并非工匠，我们却请他做一件工，他应召而来，然后动手处理备好的建筑材料，却糟蹋了木料和石材，以致房子建好随即坍塌，他能以\Quote{是别人催促我来，并非我自愿}为借口就足够了吗？绝不；而且这样做非常合理公正。因为即使他人召唤，他也应当拒绝。那么，那个只糟蹋木石的人尚且无法逃脱受罚，而那个毁坏灵魂、漫不经心建造天主圣殿的人，竟以为他人的强迫可以成为他逃避惩罚的理由？这不是极为荒谬吗？因为我还暂且不提：没有人能强迫一个不愿做的人。但即使他受到了过度的压力和种种诡计的摆布，因而落入陷阱，难道这就能使他免受惩罚吗？我恳求你们，不要自欺，假装我们不知道连小孩子都明白的事情。因为这种无知的装假在清算之日必然毫无用处。你说，你并不渴望接受这一高位，自知软弱。很好。那么，你应当以同样的心意拒绝他人的恳求；或者，当无人召唤你时，你是软弱无能的，但一旦有人愿意把这尊位给你，你突然就变得胜任了？这是多么可笑的无稽之谈！当受极刑。为此，主也劝诫那想要建塔的人，不要在没有估量自己的建造能力之前就奠基，以免给过路人留下无数讥笑他的机会。但在此处，惩罚仅仅在于成为笑柄；而在我们面前的情况，惩罚却是不灭的火、不死的虫，\BibleRef{依 66:24} 咬牙切齿、外面的黑暗、被斩割，并与伪善者同分。\par

但是控告我的人不愿考虑这些事情。否则，他们就不会责备一个不愿无端灭亡的人。眼下我们所考虑的并非谷物、大麦、牛或羊，也不是任何诸如此类的东西，而是耶稣的身体。因为按照圣保禄，基督的教会就是基督的身体；托付照管它的人，应当使它达到健康的状态和不可言喻的美，并处处留神，免得任何斑点或皱纹，\BibleRef{弗 5:27} 或其他类似的瑕疵损坏它的活力和秀丽。这不正是要使之堪当——尽人力所能——那凌驾于它之上的不朽而永受赞颂的头吗？如果那些渴望身体达到健壮状态的人需要医生和教练的帮助，以及严格的饮食、不断的锻炼和千百条其他规则（因为极小一点的疏忽就会扰乱并破坏整体），那么，那些受托照管那身体的人——该身体的斗争不是对抗血肉，而是对抗不可见的权势——若不远远超越寻常人的美德，并精通一切适合灵魂的医治，又怎能保持它健康无虞呢？\par

3. 请问，你岂不知道，那个身体比我们这肉身更容易遭受疾病和攻击，更容易朽坏，更难复原吗？为这些身体寻求医治的人，发现了各种药物，不同的器械，以及适合病人的饮食；而且，常常大气的状况本身便足以使病人康复；也有例子，适时的睡眠便省去了医生的一切劳动。但在我们面前的这个情况，却不可能考虑这些事；反之，在我们犯错之后，所指定的医治方法只有一种：就是有力地向人宣讲圣言。这就是那唯一的工具，唯一的饮食，最优良的大气。它替代药物、烧灼和切割，如果需要灼烧和截除，我们就必须使用这方法；如果这方法无效，其余一切都是枉然。藉着它，我们既唤醒沉睡的灵魂，又使它从炽热中平息；藉着它，我们削去过剩，补足缺欠，并施行一切为灵魂健康所必需的操作。至于如何妥善安排我们的日常生活，确实，他人的生活可以激起我们的效仿。但在虚假教义的事上，当有灵魂因此患病时，则极需圣言；这不仅是为了我们自己子民的得救，也是为了面对外在的敌人。诚然，若有人拥有圣神的利剑和信德的盾牌 \BibleRef{弗 6:16}，能行奇迹，并藉这些奇事堵住无耻反驳者的口，那他就很少需要圣言的帮助了；然而，即使在行奇迹的时代，圣言也绝非无用，而是必不可少的。因为圣保禄自己就使用了圣言，尽管他因奇迹而在各处如此令人惊叹；另一位属于\Quote{光荣的宗徒团体}的人也劝勉我们致力于获得这种能力，他说：\Quote{你们要常常准备好，对每一个问你们心中所怀希望的理由的人，作出答复}；而且他们众人一致把照顾穷寡妇的事托付给斯德望，唯一的原因就是他们自己可以有闲暇\Quote{为圣言的职务} \BibleRef{宗 6:4}。同样，我们也应当致力于此，除非我们确实赋有行奇迹的能力。可是，既然这种能力的迹象在我们身上丝毫未存，而四周许多敌人不断攻击我们，那么我们就必然应当用这武器武装自己，既为使自己不被敌人的箭矢所伤，也为能击伤他。\par

4. 因此，我们应该渴慕基督的圣言丰富地居住在我们内。\BibleRef{哥 3:16} 因为我们并非只为一种战斗而预备。这场战争是多种多样的，并且与种类繁多的敌人交战；并非所有敌人都使用相同的武器，也并非都采用相同的进攻方法；而必须与所有人交战的人，必须知晓所有人的诡计，同时既是弓箭手又是投石者，既是百夫长又是将军，既在行伍中又在指挥中，既徒步又骑马，既海战又攻城。在普通战争中，每个人通过执行自己承担的特定职责来击退敌人。但在此处则不然；若有人想在这场战争中取胜，他必须通晓这门艺术的一切形式，因为魔鬼深知如何通过任何一个可能未被防守的缺口引入自己的攻击者，并带走羊群。但当他看到牧者全副武装，精通神学知识，并熟知他的阴谋时，情况就不同了。因此，我们应当在各个方面严加防范：因为一座城，只要城墙环绕，就嘲笑围攻者，安居于极大的安全中；但若有人在城墙上打开一个哪怕门大小的缺口，其余城墙便无用了，尽管整个城墙都坚固如初；天主的城也是如此：只要牧者的智谋与智慧——这好比城墙——从四面保护它，敌人的一切诡计最终都陷入混乱与嘲笑，城内居民便安然无恙；但若有人能够拆毁哪怕一小部分，尽管其余部分屹立不动，毁灭也会通过那个缺口侵入全城。因为一个人若同时成了犹太人的俘虏，又何谈与希腊人激烈争辩？或者若他胜过这两者，却又落入摩尼派异端的魔掌？或者即使他证明了自己高于他们，如果那引入宿命论的人闯入并蹂躏羊群，又当如何？不必枚举魔鬼的所有异端，只要说：牧者若不精通驳斥所有异端，狼便可凭借其中任何一个进入并吞噬大部分羊群。在普通战争中，我们总是期待战场上士兵的胜利或失败。但在属灵战争中，情况截然不同。因为在那里，常常发生这样的情况：与一伙敌人的战斗却为其他从未参战、也未曾费力、一直坐着的敌人带来了胜利；而对此类事件缺乏经验的人，会犹如被自己的剑刺穿，成为朋友和敌人共同的笑柄。我将尝试通过一个例子来阐明我的意思。那些接受\Person{瓦伦提努斯}、\Person{马西昂}以及所有同样思想患病之人的狂野教义的人，将天主通过\Person{梅瑟}赐下的法律排除在圣经正典之外。但犹太人如此尊崇法律，以致尽管废除它的时期已经到来，他们仍坚持遵守其中所有内容，违背天主的旨意。然而，天主的教会避免了这两种极端，行于中道，既不被诱导置于法律的轭下，也不容忍它被诽谤，而是赞扬它——尽管它的时期已过——因为它在自己的季节尚存时是有益的。因此，要与这两种敌人作战的人，必须完全熟悉这条中道。因为若他想要教导犹太人，他们固守旧法律已不合时宜，便开始毫不留情地指责法律，他就为那些想要拆毁它的异端者提供了不小的把柄；而若他为了堵住他们的口而过度赞扬法律，并钦佩地谈论它，视其为当代所需，他便解开了犹太人的口。同样，那些患有\Person{撒伯里乌}癫狂和\Person{阿黎乌}疯狂的人，都因未持守中道而偏离了健全的信德。这两个异端都自称为基督徒；但若有人审查他们的教义，会发现一派比犹太人好不了多少，只是名称不同；另一派则几乎持有\Person{撒摩撒达的保禄}的异端，而且两者都离真理甚远。因此，在这些情形中危险极大，正教的途径狭窄，两边都有险峻的悬崖包围，并且大有恐惧，恐怕当要攻击一个敌人时，却被另一个敌人所伤。因为若有人断言天主的独一性，\Person{撒伯里乌}立即将这一表述转用于他自己的精神妄想；而若他区分位格，说圣父是一、圣子是另一位、圣神是第三位，\Person{阿黎乌}便起来，准备将这种位格的区别曲解为本体的差异；因此，我们必须转身逃离二者——逃避前者亵渎性的位格混淆，以及后者无理智的本体分割——宣认圣父、圣子、圣神的天主性是同一的，同时加上三位格的区分。这样，我们便能以抵抗两个异端的攻击。除这些之外，我可以告诉你其他一些敌人，除非我们勇敢而谨慎地作战，否则我们将满身伤痕地离开战场。\par

5. 为何要描述我们自己人的愚蠢闲谈？因为这些并不亚于来自外部的攻击，而它们给教师带来更多麻烦。有些人出于无聊的好奇，鲁莽地专注于那些既不可能知道、知道了也无益的事情。另一些人则向天主要求祂判断的说明，强迫自己探查那深不可测的深渊之底。圣经说：\BibleQuote{因为你的判断如同深渊}，你会发现很少人关心自己的信德和行为，而大多数人却好奇地探究那些无法发现、且仅探究本身便触怒天主的事情。因为我们若执意学习\Emph{祂}不愿我们知晓的事，必然失败（岂能违背天主的旨意而成功？），剩下的只有探究带来的危险。尽管情况如此，每当有人权威地阻止对这种无底深渊的探索时，他便背负上傲慢无知的恶名；因此主教此时需要极大的技巧，既要引导百姓远离这些无益的问题，又要使自己免受上述指责。简言之，应对所有这些困难，唯有言语可作帮助；若有人缺乏这能力，他所托管之人的灵魂（我指那些较软弱、好管闲事者）就像船只不断遭受风暴。因此，司铎应尽力获得这种力量。\par

6. \Person{巴西略}：\Quote{那么，为何圣保禄不渴望在这技巧上成为完美？他毫不隐瞒自己言语的贫乏，反而明确承认自己不熟练，甚至这样告诉那些以口才著称并以此自夸的格林多人。}\par

\Person{金口圣若望}：这正是败坏许多人、使他们在研习真道方面松懈的原因。因为他们未能领悟宗徒思想的深度，也未能理解他言语的含义，却终日昏睡呵欠，所尊崇的不是圣保禄所承认的那种无知，而是一种他在这世上同样完全免于的无知。\par

但暂且搁置这问题以待我们思考，我暂时只说这些。即使圣保禄在口才方面如他们所想的那样不熟练，这与今天的人有何相干？因为他拥有远超口才的能力，一种带来更大成果的能力，那就是他即使沉默，单单临在就令魔鬼恐惧。然而今天的人，即使聚集在一处，也无法用无尽的祈祷和眼泪行出圣保禄的手帕曾经行的奇迹。他藉祈祷使死人复活（\BibleRef{宗 20:10}），行了其他奇迹，以致外邦人视他为神（\BibleRef{宗 14:11}）；在他离世之前，他已配被提升到第三层天，听到凡耳不可听的言语（\BibleRef{格后 12:2}）。但今天的人——我不是要说什么严厉苛刻的话，因为我这样说并非侮辱他们，只是惊叹——当他们以此伟大人物衡量自己时，为何不战栗？因为若我们撇开奇迹，转向这位蒙福圣人的生命，察看他的天使般的交往，你会发现这位基督的斗士在这方面比在奇迹中更得胜。谁能描述他的热忱与容忍、他的不断危险、他的持续忧患、对众教会的不息挂虑、对软弱者的同情、他的诸多苦难、他不同寻常的迫害、他每日的死亡？世上何处、何大陆或海洋，不认识这义人的劳苦？连旷野也知道他的临在，因为它常在他危险时庇护他。因为他遭受了各种攻击，取得了各种胜利，他的争斗与凯旋永无止境。\par

然而，我不知不觉地竟伤害了这人。因为他的业绩超乎一切描述，尤其超乎我的能力，正如口才大师超过我。但是，既然那位圣宗徒不会以结果、而会以动机来判断我们，我就不会停止，直到我再陈述一件比之前所说的一切更甚的事，正如他自己超过所有同人一般。这是什么呢？在如此多的业绩、如此多的胜利之后，他祈祷自己可以下地狱，被交付于永罚，若能使那些曾多次用石头砸他、尽力除灭他的犹太人得救、皈依基督（\BibleRef{罗 9:3}）。有谁如此切望基督？若我们对他的情感不应描述为比切望更高尚的东西，那么我们还能再将自己与这位圣人相比吗？他领受了如此大的恩宠，并显出了如此大的德行？还有什么更傲慢的呢？\par

现在，我要在下文中竭力表明，他并非像有些人所认为的那样，是缺乏口才之人。在人眼中，缺乏口才者不仅指不谙世俗演说技巧之人，也包括无力为正确信仰辩护之人，而他们是对的。但圣保禄并未说他在这两方面都缺乏口才，而仅在一方面；为此，他仔细区分，说自己\Quote{言语粗俗，然而知识却不粗俗。}\BibleRef{格后 11:6} 倘若我坚持任何一位主教必须拥有伊索克拉底的文雅、狄摩西尼的雄浑、修昔底德的高贵和柏拉图的深邃，那么圣保禄便是驳斥我的有力证据。但我撇开这一切世俗演说的精致修饰；我并不计较文体或口才；是的，一个人的措辞可以贫乏，文笔可以简朴无华，但不可在知识的准确表述上无知；也不可为了掩饰自己的懈怠，而剥夺那位圣宗徒最伟大的恩赐和全部赞颂之精华。\par

7. 请告诉我，他尚未开始行奇迹，又是如何使住在大马士革的犹太人惊惶失措？\BibleRef{宗 9:22} 他如何与希腊人辩论并胜过他们？他为何被派往塔尔索？岂不正是因为他言辞如此有力而胜利，使他的对手无法忍受失败，竟要谋害他的性命？当时，如我所说，他尚未开始行奇迹，无人能说群众因他行大能的荣耀而惊叹，或是与他辩论的人被他公众声望的力量所压倒。因为此时他仅仅靠论证取胜。再者，他如何与那些想在安提约基雅推行犹太化的人成功辩论？那位雅典的亚略巴古议员\BibleRef{宗 17:34}——那座最虔诚于神明的城市——和他的妻子又是如何跟随了宗徒？岂不因他们所听的讲论？当厄乌提科\BibleRef{宗 20:9}从窗口坠下时，岂不因他听保禄讲道直到半夜？我们在得撒洛尼、格林多、厄弗所，乃至罗马本地，见到他在做什么？难道他不是整日整夜地按次序解释圣经？至于他与厄壁鸠鲁和斯多葛派的辩论\BibleRef{宗 17:18}，又何必详述？若我们决意细述每一细节，故事将变得冗长不堪。\par

因此，当圣保禄在行奇迹前后都大量运用论证时，谁还能胆敢断言他缺乏口才？他的讲论和辩论如此深得所有听众的钦佩。何以吕考尼雅人\BibleRef{宗 14:11}会以为他是赫尔默斯？他们以为他和巴尔纳伯是神，确实出于见到奇迹；但以为他是赫尔默斯，并非因此，而是出于他的言语。这位蒙福的圣人在哪方面超越其他宗徒？何以他在普世被交口称赞？不仅在我们中间，在犹太人和希腊人中，他为何是奇迹中的奇迹？岂不来自他书信的力量？藉此，不仅对今日的信友，而且从那时直到如今，直到基督的显现，他过去、现在和将来都有益于人，并将继续如此，只要人类存在。因为他的著作犹如金刚石筑成的墙，坚固了普世的教会；他如同最尊贵的勇士，站立其中，将一切思想夺回，使之顺服基督，攻破诡辩，以及一切高抬自己、反对天主之知识的堡垒\BibleRef{格后 10:5}，这一切他都藉着留给我们的书信完成，这些书信充满了奇迹和属天的智慧。他的著作不仅有益于我们推翻错误教义、确认真道，而且在度圣善生活方面也大有帮助。如今的主教们利用这些书信，塑造那位贞洁的童女——就是圣保禄亲自许配给基督的\BibleRef{格后 11:2}——引导她达致属灵的美丽状态；并藉此驱散围绕她的有害瘟疫，保持所获的健康。那位被视为缺乏口才的人留给我们的药物及其效能，正是经常使用它们的人才深知其功用。由此可见，圣保禄曾以极大的勤勉和热忱致力于我们所说之研究。\par

8. 请听他如何吩咐他的门徒：\BibleRef{弟前 4:13} \Quote{你要专注于宣读、劝勉、教导；}他又接着指出这样行的益处，附加说：\Quote{因为这样做，你不但能救自己，也能救那些听你的人。}\BibleRef{弟前 4:16} 他又说：\Quote{主的仆人不应争斗，而应温良对待众人，善于教导，忍耐；}\BibleRef{弟后 2:24} 他继续说：\Quote{但你应坚持你所学的、所确信的，你知道你是从谁学的，并且自幼认识圣经，这圣经能使你有得救的智慧。}\BibleRef{弟后 3:14} 又说：\Quote{凡圣经都是天主所默感的，对于教训、督责、矫正、教导人学义，都是有益的，为使天主的人得以完备。} 请听他在致弟铎书关于选立主教的指示中又加上了什么。他说：\Quote{主教必须持守那合乎教导的忠信之言，好能驳斥反对的人。} 但假如一个人如同这些人所假装的那样不熟练，又怎能驳斥反对的人并堵住他们的口呢？或者，如果这种不熟练的状态在我们中间受欢迎，那么专心阅读和研读圣经又有什么必要呢？这些论点不过是权宜之计和借口，是懒惰和懈怠的标志。但有人会说：\Quote{这些吩咐是给司铎的。}——的确，因为他们是我们要讨论的对象。但是，宗徒也将同样的吩咐给了平信徒，请听他在另一封致非司铎群体的书信中说什么：\Quote{要让基督的话丰丰富富地存在你们心里，以各样的智慧；}\BibleRef{哥 3:16} 又说：\Quote{你们的言语要常常带着恩宠，用盐调和，使你们知道应当怎样回答各人。}\BibleRef{哥 4:6} 并且有一个普遍的吩咐，要众人\Quote{随时准备}\BibleRef{伯前 3:15}为自己的信仰作答；又给得撒洛尼人以下命令：\Quote{你们要彼此劝慰，互相建立，正如你们一向所行的。}\BibleRef{得前 5:11} 但当他论及司铎时，他说：\Quote{那些善于管理的长老，尤其是那些在言语和教导上劳苦的，应被视为配享双倍的尊敬。}\BibleRef{弟前 5:17} 因为当教师既凭行为又凭言语带领门徒达到基督为他们所安排的福境时，这就是教导的完美境界。因为仅仅榜样不足以教导他人。我这样说不是出于自己，而是出自我们救主的亲口之言。因为\Quote{谁若实行并教导，他将被称为大。}\BibleRef{玛 5:19} 如果行为等同于教导，那么这里第二个词就是多余的；只要简单说\Quote{谁若实行}就够了。然而，他区分了二者，表明实践是一回事，道理是另一回事，且两者需要彼此帮助以完成全备的建立。你也听到基督特选的器皿向厄弗所的长老们说了什么：\Quote{所以你们应当警醒，记住我三年之久，昼夜不停地流泪劝诫你们每个人。}\BibleRef{宗 20:31} 既然他作为宗徒的生活如此光辉，那么他的流泪和口头劝诫还有什么必要呢？他的圣德生活也许能极大地激励人遵守诫命，然而我不敢说单靠它就能成就一切。\par

9. 但当有关教义的争论发生时，所有人都从同一圣经中取用兵器，一个人的生活能证明什么分量呢？那么，他许多克苦的修行又有何用，当经过如此艰苦的训练之后，任何人因司铎在辩论上极不熟练而陷入异端，并从教会身体中被切断——这是我自己曾见许多人遭受的不幸。那时他的忍耐对他有何益处呢？毫无益处；正如如果生活腐败，纯正的信仰也无益一样。因此，比任何其他理由更甚，那负有教导他人职责的人，应当精通此类辩论。因为即使他本人站得稳，没有被反对者伤害，然而他领导下的单纯群众，当他们看到自己的领袖被击败，对反对者无言以对时，便容易将他失败的责任归咎于道理本身，好像道理有缺陷似的；于是因一人的不熟练，很多人陷入极端的毁灭；因为他们虽未完全投靠敌人，却被迫对自己从前坚信不疑的事产生怀疑，而那些他们曾以坚定信心接近的事，再也无法坚定地持守，反而因领袖的失败，一场巨大的风暴降临到他们的灵魂上，以致这祸患最终使他们完全沉船。但是，这样的领袖为他所丧失的每一个灵魂，给自己可怜的头颅带来何等可怕的毁灭，何等可怕的火，你不需要从我这里学习，因为你完全知道这一切。那么，我不愿成为这么多灵魂灭亡的原因，并为自己在来世招致比现在已等候我的更重的惩罚，这难道是骄傲，是虚荣吗？谁会这样说呢？当然没人会，除非他希望在无错之处挑剔，并对他人的灾难加以说教。\par

\SourceRecord{Translated by W.R.W. Stephens.；Nicene and Post-Nicene Fathers, First Series；Vol. 9.；Edited by Philip Schaff.；Buffalo, NY: Christian Literature Publishing Co.,；1889.；<http://www.newadvent.org/fathers/19224.htm>.}

% structure: p0001 source=h1 target=section reason=page-title

\section{论司祭职（第五卷）}

1. 上文已充分说明了教师为真理热切辩论需要多么大的技巧。但我还必须提及另一件事，这件事是无数危险的根源，不过就我而言，我宁愿说事情本身并非原因，而是那些不知如何正确使用它的人，因为在那些热忱而善良的人管理下，它本身就是救恩和许多其他善事的帮助。那么，我指的是什么呢？就是在公共场合发表讲道需要付出大量的辛劳。首先，大多数受讲道者照管的人并不愿意以对待教师的态度对待他们，反而摒弃学习者的角色，取而代之的是那些坐着观看公共比赛的人的态度；正如那里的群众分成党派，有人附和这个人，有人附和那个人；同样，这里的人也分裂了，成为此时这位教师、那位教师的拥护者，带着好感或恶意听他们讲。不仅如此，还有另一种几乎同样严重的困难。因为如果有哪位讲道者将别人的作品部分编入自己的讲道中，他比那些偷钱的人遭受更大的耻辱。不，即使他没有向任何人借用任何东西，只是被怀疑，他也会遭受窃贼的命运。我为什么还要提别人的作品呢？连他自己也不能没有变化地使用自己的资源。因为公众习惯了不是为益处而听，而是为乐趣而听，像悲剧和音乐表演的批评家一样坐在那里；我们刚才痛斥的那种口才，在这种情况下变得比律师更受欢迎，在律师那里他们必须彼此争辩。讲道者因此应该有高尚的心志，远超过我卑微的精神，使他能够纠正群众身上这种混乱而无益的乐趣，并能够引导他们转向更有益的听道方式，好使他的百姓跟随并顺从他，而不被他自己的情绪牵引；而除了两种方法外，不可能达到：漠视他们的称赞，以及擅长讲道的才能。\par

2. 如果缺乏这两者之一，剩下的一个就会因与另一个分离而变得无用；因为如果讲道者漠视赞扬，却不能产生那\BibleQuote{带着恩宠调上盐，}\BibleRef{哥 4:6}他就会受到群众的鄙视，同时从他的高尚心志中一无所获；另一方面，如果他是一位成功的讲道者，但被掌声的想法压倒，伤害也同样发生，既对他自己，也对群众，因为他渴望赞扬，便谨慎地说话，与其说为了利益，不如说为了取悦。正如那既不受好评影响，又不擅长讲话的人，不屈服于群众的乐趣，也无法给他们带来任何值得一提的好处，因为他无话可说；同样，那个被赞扬欲望冲昏头脑的人，虽然能够更好地服务群众，却反而提供适合他们口味的食物，因为他借此购买喝彩的喧闹。\par

3. 因此，最好的主教必须在这两点上都坚强，使二者互不相害。因为如果当他站在会众中，说出旨在使疏忽的人畏缩的话，却随即结巴、停顿，因失败而被迫脸红，那么他所说的益处就立刻浪费了。因为那些受责备的人，被所说的话刺痛，又无法以其他方式向他报复，就嘲笑他的无知，企图借此掩盖自己的耻辱。因此，他应当像一位极其优秀的车夫，对这两种善事作出精确的判断，以便按需处理二者；因为当他在众人眼中无可指责时，他就能随意使用他的权威，纠正或免除纠正所有在他治下的人。但没有这一点，他就不容易做到。而这种灵魂的高尚不仅应当表现在漠视赞扬的程度上，而且应当更进一步，使这样获得的益处不至于再次变得无果。\par

4. 他当对什么事漠不关心呢？毁谤与嫉妒。然而，对于不合时宜的诽谤（因为主教难免会遭受无端的指责），他既不应过分畏惧和颤抖，也不应完全忽视；我们应当努力立即扑灭它，即使它是虚假的，或是由普通民众加给我们的。因为没有什么比不受管束的民众更能夸大恶名和美名了。他们习惯于不经调查就听和说，随意重复一切遇到的事，毫不顾及事实。因此，主教不应当对群众漠不关心，而应当立即将他们的恶意猜测扼杀在萌芽状态，说服他的控告者，即使他们是世上最不讲理的人，并且不应忽略任何能够驱散不利传闻的事。但当我们做了这一切后，责备我们的人仍不肯被说服，那么从那时起，我们就不应再关心他们。因为如果有人因这些意外而过于沮丧，他永远不能产生任何崇高和可钦佩的事。因为消沉和不断的忧虑对摧毁心智的力量并将其带入极端软弱极有帮助。因此，司铎对待他管辖的人，应当像父亲对待幼小的子女一样；他们不会因子女的辱骂、击打或哀哭而烦恼，甚至当子女与我们一同欢笑和喜悦时，我们也不十分在意；同样，我们既不应因这些人的许诺而自高，也不应因他们的指责而沮丧，即使这些指责是不合时宜的。但这很难，我的好友；而且，我想，甚至是不可能的。因为我不知道是否有人曾成功做到在被赞美时而不喜悦，而为此喜悦的人也容易渴望享受这份喜悦，渴望享受它的人必然会在失去它时感到烦恼和失魂落魄。因为正如那些以富有为乐的人，当他们陷入贫穷时会悲伤；那些习惯于奢华生活的人无法忍受简陋的生活；同样，渴望喝彩的人，不仅在被无故指责时，而且在没有持续被称赞时，灵魂就会像遭受饥荒一样消瘦，特别是当他们自己习惯于被赞美，或当他们听到别人被赞美时。怀着这种渴望开始讲道试炼的人，你认为他有多少烦恼和多少痛苦呢？海洋不可能没有波浪，那人也不可能没有挂虑和忧伤。\par

5. 因为即使讲道者可能拥有极大的能力（这只能在少数人身上找到），即使这样，他也不能免于不断的辛劳。因为讲道并非出于天生，而是由于学习，假定一个人达到了很高的水准，那么如果他不通过不断的运用和操练来培养自己的能力，这水准就会离他而去。因此，智慧者比无知者承受更大的辛劳。因为疏忽对两者的损失程度并不相同，损失与两者所拥有的差异完全成正比。对于后者，没有人会责备，因为他们提供不出值得重视的东西。但对于前者，除非他们不断产出超越所有人对他们声誉期望的内容，否则他们就会受到各方的严厉指责；此外，后者即使是小小的表现也会获得相当的赞扬，而前者的努力，除非特别精彩和惊人，不仅得不到掌声，反而遭遇许多挑剔者。因为听众以批评者自居，与其说是评判所说的内容，不如说是评判讲道者的声誉；因此，当有人演讲能力超过众人时，他尤其比其他人都更需要刻苦的学习。因为这个人不允许利用人性常有的借口——人不能事事成功；但如果他的讲道辞整个不符合所形成期望的巨大程度，他将一无所获地离开，只得到无数的嘲笑和指责；没有人会考虑到这一点：沮丧、痛苦、焦虑，甚至愤怒，可能插入进来，使他的思想清晰度变暗，并阻止他的作品以纯粹的形式从他那里出来；而且总的说来，他只是一个人，不可能始终如一，也不可能一直都成功，而是必然有时会未达目标，表现出比平时低的能力水平。正如我所说的，他们不愿考虑这些，却指责他的过错，好像他们坐在审判一位天使；尽管在其他情况下，人也容易忽略邻人的出色表现，即使这些表现又多又好；而一旦出现一个缺陷，即使是偶然的，即使只是很久才出现一次，也会很快被察觉并永远被记住；因此，细小琐碎的事往往减损了许多伟大业绩的荣耀。\par

6. 你看，我卓越的朋友，那在宣讲上有大能的人，特别需要比别人更大的勤学；除勤学外，也需要比以上我所提及的一切人更大的容忍。因为常有许多人，以虚空无知的精神起来攻击他，找不出他的过失，只因他普遍受人称赞，便恨他；他必须勇敢地忍受他们恶毒的怨恨，因为他们既不能隐藏这无理怀有的咒诅般的仇恨，便私下毁谤、责难、中伤，公开诋毁。那开始因每一次这样的事而痛苦和恼怒的心灵，必将因忧伤而败坏。因为他们不仅用自己的行为报复他，还要假手他人来报复；常常选出一个不善于言辞的人，用赞美将他捧得过高，超出他的价值。有人纯因无知而行，有人因无知兼嫉妒而行，为的是毁坏那另一人的声誉，并非要证明那人其实并不出色却显得奇妙；而心胸高尚的人不仅要与这些人斗争，还要常与全体听众的无知斗争。因为来聚会的人不可能都是博学之士，多数情况是大部分会众由无学问的人组成，即使其余头脑较清晰的人，与前者相比也远不能批评讲道，而其余的人又不如他们；所以只有一两个人坐在那里具有这种能力；由此必然，那讲得比别人更好的人所得的掌声较少，甚至可能回家时全无称赞；他必须勇敢地应对这种反常，宽恕那些因无知而犯错的人，为那些因嫉妒而默认的人感到可怜可悯，而不因这两种原因之一认为自己的能力减弱了。因为如果一个人是卓越的画师，在技艺上超越众人，他精确绘制的肖像被无知者讥笑，他不应沮丧，认为画作差劲，因无知者的判断；正如他不会因无艺术修养者的惊异，就把真正拙劣的画作当作奇妙可爱的。\par

7. 让最优秀的工匠成为自己设计的评判者，让他的作品根据设计它的心智所作的判断被定为优劣。甚至不要考虑外界如此错误而无艺术修养的意见。因此，承担教导辛劳的人，切不可在意外界的好评，也不要因这样的人而心中沮丧；要致力于他的讲道，为能悦乐天主（因为让这作为他履行这最佳工艺的规则和定断，而非喝彩或好评），诚然，若被人称赞，他不应拒绝掌声；而听众不提供时，他不应寻求，也不应忧伤。因为在他的劳苦中，有一个充分的、超越一切的安慰，就是能意识到自己的教导是为悦乐天主而安排和秩序井然的。\par

8. 因为如果一个人首先被对无差别赞美的渴望所冲昏头脑，他就不会从他的劳苦或他讲道的权能中获得任何益处，因为心灵无法忍受群众无理的指责而沮丧，并将所有关于讲道的热忱抛诸脑后。因此，特别需要训练自己对各种赞美都漠不关心。因为知道如何讲道，对于保持那种权能来说是不够的，除非再加上这一点：如果有人要准确审视那缺乏这技巧的人，他会发现，他需要不亚于前一个人那样对赞美漠不关心，因为他若将自己置于舆论的控制之下，就会被逼做出许多错事。因为若没有能力与那些在讲道品质上享有声誉者匹敌，他就会忍不住对他们怀有恶意、嫉妒他们、无端指责他们，以及许多此类不光彩的行为，并且会不惜一切，甚至需要毁灭自己的灵魂，也要将他们的名声拉低到自己卑微的水平。除此之外，他会停止对自己工作的努力；一种麻木，可以说，蔓延到他的心灵。因为辛苦劳作却得到微薄的赞美，足以击倒一个不能轻视赞美的人，使他陷入深深的昏睡；正如农夫在花费时间耕种一块贫瘠的土地，被迫犁开岩石时，除非他对此事怀有极大的热忱，或担心饥荒逼近，否则很快就会放弃劳作。因为那些能相当有力讲话的人，尚且需要如此不断的锻炼来保持他们的才能；而一个根本不收集材料，却被迫在努力中思考的人，要经历多大的困难、混乱和麻烦，才能靠巨大努力收集到一点点想法！如果在属下的神职人员中，有人虽处于较低阶位，却能在那个位置上显得比他有更好表现，那么他需要拥有何等神圣的心智，才能不被嫉妒攫住或因沮丧而跌倒。因为一个人被置于更高的尊位，却被地位较低的人超越，并能高贵地承受这一点，这不是任何普通心智所能做到的，也不是像我这样的心智所能做到的，而是需要像金刚石一样坚硬的心智。而且，如果那位更著名的人非常宽容和谦逊，那么这种痛苦就变得更容易承受。但如果他大胆、自夸、虚荣，那么对于另一个人来说，每天死去都是可取的；他将如此毒害他的生活，当面侮辱他，背后嘲笑他，从他那里夺走许多权力，并希望自己成为一切。但在所有这些情况下，拥有最稳妥保障的人，是那些在讲道上有流利口才、得到群众热切关注、并受到所有属下爱戴的人。难道你不知道如今一种对讲道的热情已经涌入基督徒的心中吗？那些练习讲道的人得到特别的尊荣，不仅在外邦人中，也在信德之家的人中。那么，如果有人发现自己讲道时所有人都沉默，认为他们受压迫，等待讲道结束如同等待从工作中解脱；而听别人讲道时却热切渴望，即使对方讲道很长，他们也为他将要结束而感到遗憾，并且几乎对他打算沉默感到愤怒，他怎能忍受这样的羞辱呢？如果这些事情在你看来微不足道，容易轻视，那是因为你缺乏经验。它们确实足以冷却热忱，麻痹心灵的力量，除非一个人远离一切人类的情感，并努力按照那些无形体之能的模样塑造自己的行为，它们既不受嫉妒追逐，也不受对名声的渴望或任何其他病态情感的影响。因此，如果有人能这样被造就，能够制服这难以捕捉、不可征服、凶猛的野兽——也就是说，公众名声——并砍下它的许多头，或者更确切地说，完全阻止它们生长，他将能够轻易地击退这些许多猛烈的攻击，并享受一种平静的安息港湾。但是，凡没有从这怪物中解放出来的人，就会使他的灵魂陷入各种斗争、不断的动荡以及沮丧和其他情欲的重担之中。但我何必详细叙述其余的困难呢？这些困难没有人能够描述，除非他亲身经历过。\par

\SourceRecord{Translated by W.R.W. Stephens.；Nicene and Post-Nicene Fathers, First Series；Vol. 9.；Edited by Philip Schaff.；Buffalo, NY: Christian Literature Publishing Co.,；1889.；<http://www.newadvent.org/fathers/19225.htm>.}

% structure: p0001 source=h1 target=section reason=page-title

\section{论司铎职（第六卷）}

我们的处境确实如你所闻。但将来的处境，当我们被迫为每一个托付给我们的人交账时，我们怎能承受？因为我们的惩罚不仅限于羞耻，永罚也在等待我们。至于那段经文：\BibleQuote{你们应信服并服从那些带领你们的人，因为他们为你们的灵魂时刻警醒，好像要交账的人一样} \BibleRef{希 13:17} 虽然我已提过一次，但现在我仍要打破沉默，因为那警告的恐惧不断搅动我的灵魂。因为若有人使一个最小的信者跌倒，倒不如\BibleQuote{把大磨石拴在他颈上，沉在深海里} \BibleRef{玛 18:6}；且若有人伤及弟兄的良心，就是得罪基督自己（\BibleRef{格前 8:12}），那么，那些毁掉的不只是一人、两人或三人，而是如此众多人群的人，有朝一日将受何苦，将付何代价？因为不能以经验不足为借口，也不能以无知为托辞，也不能提出必要或强制的理由。是的，如果可能，在他们管辖下的某个人倒更容易为自己的罪过使用这种托辞，而不是主教为别人的罪过。你问为什么？因为被任命来纠正他人无知、预先警告他们即将到来与魔鬼的争战的人，将不能以无知为借口，也不能说：\Quote{我从未听见号声，我没有预见争战}。他正是为此被设立的——\Person{厄则克耳}说——好为他人吹号，警告他们眼前的危险。因此他的惩罚是不可避免的，即使迷失的只有一人。\BibleQuote{当刀剑来时，守望者若不向人民吹号，也不给予信号，刀剑就来把人卷去，他固然因自己的罪而被卷去，但我要向守望者追讨他的血} \BibleRef{则 33:6}。\par

那么，请不要催促我们走向这样必然的惩罚；因为我们所说的并非军队或王国，而是一项需要天使般品德的职务。因为司铎的灵魂应当比太阳光还要纯洁，好让圣神不离开他，使他能说：\BibleQuote{我如今活着，但不再是我，而是基督在我内活着} \BibleRef{迦 2:20} 因为如果那些住在旷野、远离城市和市场及其喧嚣、始终享有安息与平安港湾的人，尚且不满足于那种生活方式的安全，反而加上无数其他防护，四面环绕自己，并研究以极其谨慎的方式说话和行事，以便尽人之所能，怀着确据和无玷的纯洁亲近天主，那么，你认为被祝圣的司铎需要怎样的力量和坚韧，才能将自己的灵魂从一切污秽中撕裂，并保持其属灵之美不受玷污？因为他需要的纯洁远远超过他们；谁需要更大的纯洁，谁也就面临比他们更紧迫的诱惑，这些诱惑能玷污他，除非他通过不断的克己和大量劳苦，使自己的灵魂对它们不可接近。因为容貌的美、步态的优雅、做作的步态和嗲声、画眉与涂颊、精心的编发和染发、衣着的昂贵、金饰的多样、宝石的光彩、香水的芬芳，以及妇女们所投身的一切其他事物，都足以扰乱心神，除非它碰巧因极其严格的自我克制而变得坚强。被这些事物所扰乱并不奇怪。但另一方面，魔鬼能够通过这些的对立面击中并射倒人的灵魂——这却令我们感到惊讶和困惑。\par

因为此前有些人逃脱了这些罗网，却被其他截然不同的罗网所捕获。即使是邋遢的外表、不梳理的头发、肮脏的衣服、未施脂粉的脸、简单的举止、朴实的言语、不刻意造作的步伐、不矫揉造作的嗓音、贫穷的生活、卑微、无人问津和孤独的境况，也首先吸引旁观者产生怜悯，进而导致彻底的毁灭；许多逃脱了金饰、香水、衣着等前述罗网的人，却轻易落入这些截然不同的罗网中而丧亡。既然如此，当贫穷与富足、装饰与不修边幅、刻意与自然，总之我所列举的一切，都在观看者灵魂中引发战争，诡计从四面环绕他时，他怎能自由呼吸，当如此多的罗网包围他时？他又能找到怎样的藏身之处——我不是说避免被强行捕获（这并非完全困难），而是说使自己的灵魂不受污秽思绪的干扰？\par

再者，我略过种种尊荣不提，那些尊荣正是无数祸患的根源。由妇人之手而来的尊荣，有损于自制之力，且常能将其倾覆，倘若一个人不知常加防范此类计谋；而由男人之手而来的尊荣，除非人以十分高贵的胸怀去接受，否则他会被两种相反的情绪——奴颜婢膝的谄媚与愚昧的骄傲——所攫取。他向庇护他的人不得不卑躬屈节，而因他人所给予的尊荣，便对属下趾高气扬，跌入傲慢的深渊。我们固然提及了这些事，但它们实际有多么有害，若非亲历，无人能清楚知晓。因为一个人若在世上生活，所必须面对的不仅这些陷阱，更有比这些更大、更具迷惑性的圈套。然而，安于独处的人，却能摆脱这一切；即便偶尔有怪异的念头在心中生出这类的景象，那影像也是微弱的，且能迅速被制服，因为外在的实际所见并未给火焰添加燃料。因为隐修者只需为自己担忧；倘若他被迫照料他人，那些人也屈指可数；即或人数众多，也比我们教会中的人少，且为他们的长上所操的心也轻省得多——这不仅因为他们人数少，更因为众人都摆脱了俗务，没有妻子、儿女或任何此类牵挂；这使他们对自己的长上极为顺从，并得以同住一处，因此能一目了然地察觉并纠正他们的过失，因为教师持续的监督对于德行的进步并非小助。\par

4. 然而，隶属于司铎的那些人，大多为今生的挂虑所累，这使得他们在履行属灵本分时较为迟缓。因此，教师必须每日（可以这么说）播种，以便藉着频繁的宣讲，至少能使听道者把握住教训之言。因为过度的财富、充裕的权势、奢靡所生的怠惰，以及除此以外的许多事物，都扼杀已撒落的种子。荆棘也常常长得过密，连种子落地都不得。此外，过度的烦恼、贫困的压力、不断的侮辱，以及其他与此相反的事，又将人的心思从对神圣之事的挂虑上夺走；而民众的罪过，即便是最小的一部分，也难以显明——因为在那些大多数连面目都不认识的人身上，又怎能察觉呢？\par

司铎与民众的关联已包含如此多的困难。但若有人探究他与天主的关联，便会发现前述种种都不值一提，因为后者需要更重大、更彻底的认真。因为那为全城——我为何说为全城？实则为全世界——充当大使的人，祈求天主垂怜所有人的罪过，不仅活着的人，也包括已亡者。他该是怎样的人呢？依我看来，连\Person{梅瑟}和\Person{厄里亚}的坦率直言，对于这样的恳求也嫌不足。因为他仿佛受托照料全世界，且是众人的父亲，走近天主，恳求战火随处熄灭，骚乱得以平息；求赐平安与丰裕，以及一切临到众人、公开或私下之灾祸的迅速解救；他必须在各方面远超过他所代祷的众人，正如统治者应超过其臣民一般。\par

每当他呼求圣神，举行最可敬畏的祭献，并不断接触万有的共同主时，请问我们该给他什么地位呢？我们必须要求他有何等纯洁、何等真正的虔诚呢？试想，从事这些事的手应当是怎样的一双手，说出如此话语的舌头又当是怎样的舌头，那领受如此伟大圣神的灵魂，难道不应当比世上任何事物都更纯洁、更神圣吗？在这样的时刻，天使侍立在司铎身旁；整个圣所和祭台周围充满了天上的军旅，以尊崇那躺在祭台之上的那一位。这确实可以从当时举行的礼仪本身得到证明。我自己也曾听人讲述，一位可敬的长者，惯常看到神视，曾告诉他说，他自己有幸得见这一类的异象，那时，他突然看到——尽可能地——一大群天使，穿着闪亮的衣袍，环绕着祭台，屈身下拜，如同人们看到士兵在君王面前那样。我对此深信不疑。另一个人也告诉我，他不是从别人那里听说的，而是他自己有幸亲耳听闻、亲眼目睹，那即将离世的人，倘若以纯洁的良心领受了圣事，在他们将要咽气时，天使因他们所领受的而守护他们，并护送他们离去。你难道还不战栗吗？竟要将一个穿着污秽衣裳、被基督从宾客行列中赶出去的人，引入如此神圣的奥迹，并提拔到司铎的尊位？见：\BibleRef{玛 22:13} 司铎的灵魂应当如照亮世界的光。但我的灵魂却因邪恶的良心而笼罩在巨大的黑暗中，以至于总是沮丧，不敢坦然仰望它的主。司铎是地上的盐。见：\BibleRef{玛 5:13} 但除了你——素来超乎寻常地爱我的你——谁又能轻易容忍我的愚昧和万事上的无经验呢？因为司铎不仅应当如此纯洁，如同身负如此崇高职务的人，而且应当非常审慎，精通许多事务，既像世间人那样熟悉世事，却又比独居山中的隐修士更远离这一切。因为他必须与那些有妻子、养育儿女、拥有仆人、富有资财、身居要职、有权有势的人交往，所以他本人也应当是一个多面手——我说多面手，并不是说虚伪、谄媚、假意，而是充满极大的自由和自信，知道如何在情况需要时有益地适应，既温和又严厉，因为不可能用同一种方式对待所有受自己管辖的人，正如医生不能对所有病人用同一种疗法，舵手也不能只知道一种对付风浪的方法。事实上，这艘船不断遭到风暴，这些风暴不仅来自外部，也来自内部，需要极大的迁就和谨慎，而这些不同的事情都有一个共同的目的：天主的荣耀和教会的建立。\par

隐修士们进行的战斗是巨大的，他们的劳苦也很多。但若有人将他们的努力与司铎职的恰当行使相比，他会发现其间的差异如同君王与平民之间的距离。因为在前者中，虽然劳苦确实巨大，但战斗是灵魂和肉体共同的，或者更确切地说，大部分是通过身体的状况完成的；如果身体不坚强，意愿便无法发展，不能付诸行动。因为严厉的禁食、睡在地上、守夜、不沐浴、巨大的劳动以及他们为克苦身体而使用的其他一切方法，在身体不够强壮时都会放弃。但在此处，所涉及的是灵魂的纯洁，并无需身体的力量来显示其卓越。因为身体的力量对使我们不自私、不骄傲、不固执、而是清醒、审慎、有序，以及 \Person{圣保禄} 所描绘的完美司铎的其他一切品德有何帮助呢？但对于隐修士的德行，没有人能这样说。\par

如同在表演奇术的人那里，需要许多器械，如轮子、绳索和匕首；而哲学家则将其全部艺术储存在心中，不需要任何外在器具。同样，在我们所讨论的情况中也是如此。隐修士需要良好的身体状况和适合其生活方式的地点，以便他们不与世人交往过远，享有旷野的宁静，此外还要有最适宜的气候。因为对于因禁食而衰弱的身体，没有比气候变化不定的气候更难忍受的了。至于他们在准备衣物和日常饮食上不得不承受的麻烦，因为他们自己渴望亲手做一切，我现在不必多言。但司铎不需要这些来供给自己的需要，他对它们毫不在意，并参与一切无害的事，同时将他所有的技艺储存在心中的宝库里。若有人赞美独处和远离人群的生活，我本人会说，这种生活是忍耐的标志，但不是完全灵魂刚强的充分证明。因为那个在港口掌舵的人，还不能给出其技艺的确定证据。但若有人能在海中安全地引导他的船，没有人会否认他是一位优秀的舵手。\par

7. 因此，我们绝不至于过分惊讶，独修者独自生活，不受搅扰，不犯众多大罪。因为他不遭遇那些刺激和激动心灵的事物。但若有人将自己投身于众人之中，被迫承担许多人的罪过，却仍坚定不移，在风暴中如同在风平浪静时一样引导自己的灵魂，这样的人理应从众人那里得到公正的赞扬和钦佩，因为他已充分证明了自己的刚毅。所以，你也不必奇怪，我既然避开市场和人群，就没有多少人指责我。因为我本不该奇怪，若我在睡眠中未犯罪，未摔跤时未跌倒，未战斗时未受伤。因为，请问，谁能起来反对我，揭露我的败落？是这个屋顶或小室吗？它们不能说话。是我的母亲吗？她最了解我的事情。然而，我与她既不通信，也从未争吵过；即使如此，也没有哪位母亲如此冷酷、缺乏对孩子的慈爱，以致在众人面前辱骂和指控她所怀孕、生产、养育的孩子——除非有理由迫使她，或有人怂恿她这样做。尽管如此，若有人愿仔细审视我的心灵，他必会发现其中有许多腐败之处。你对此并非不知，你常在众人面前特别颂扬我。我说这些并非出于谦逊，你当记得，在我们讨论此事时，我曾多次对你说：若有人让我选择，是愿意在管理教会中得荣耀，还是在独修生活中得荣耀，我会一千次投票选择前者。因为我从未停止祝贺那些能够善尽此职的人，没有人能否认，如果我能够相称地参与其中，我不会回避我所认为有福之事。但我能做什么呢？没有什么比我的这种懒散和疏忽更有损于管理教会了——别人以为这是某种自律，而我认为它不过是我个人软弱的遮羞布，掩盖了我大部分的缺陷，不让它们显露出来。因为一个惯于享受如此大的闲暇、在极大安宁中度过时光的人，即使天性高贵，也会因缺乏经验而困惑不安，他的懒散剥夺了他不少天生的能力。但若他再加上迟钝的智力，且对这些严峻的考验一无所知——我认为我就是这种情况——那么他履行所受的牧职，就不过像一尊雕像。因此，那些从这种学校出来、面临如此重大考验的人，鲜有闪耀者；大部分人都暴露了自己，跌倒，经历许多艰辛和痛苦；这并不奇怪。因为考验和训练并不涉及相同的事物。那正在争战的人，与未经训练的人毫无区别。如此进入竞技场的人，应当轻视荣耀，超越愤怒，充满极大的明智。但对于追求独修生活的人，这些品质没有施展的余地：他没有许多激怒他的人，以便练习克制愤怒的力量；也没有许多崇拜者和赞扬者，以便训练自己轻视群众的称赞。至于教会在明智方面的要求，在他们身上无从考量。因此，当他们遇到从未实际经历过的考验时，就感到困惑，头晕目眩，陷入无助的状态；不但没有增加任何卓越之处，反而常常失去了自己原有的东西。\par

8. \Person{巴西略}：那么，我们是否该将管理教会之职交给那些活跃于社会、关心世事、擅长争辩和辱骂、满腹诡计、精通奢华之道的人呢？\par

金口圣若望：安静，我亲爱的朋友！当讨论司铎职时，你决不应在心中怀有此类人，而只应怀有那些能与所有人交往、相处之后，仍能保持其纯洁无损，其超脱、圣洁、恒心和清醒不动摇，并比隐修士更甚地拥有所有属于隐修士的其他德行的人。他有许多缺陷，但能藉着隐居将其隐藏，并使其无效，因为他不同任何人交往；当\Emph{他}进入社会时，他将一无所获，只会成为笑柄，还会冒更大的风险——我几乎亲身经历了这些，若不是天主的眷顾迅速将这样的火从我头上移开。因为处在这种地位的人不可能不被察觉，当他如此显眼地处于高位时，一切都会被识破；正如火检验金属的材料，圣职的考验也探寻世人的灵魂；如果有人易怒、卑鄙或贪图名声，如果他自夸或有其他类似缺点，圣职会将一切揭示出来，并迅速暴露其缺陷，不仅暴露，而且增加其痛苦和强度。因为身体的伤口如果被擦伤，就变得更难愈合；而心灵的激动在被激怒和惹恼时，自然更加愤怒，拥有它们的人被驱使犯下更大的罪。因为那些不加以约束的人会导致他们爱慕虚荣、自夸、贪恋世俗财物，并把他们往下拉，拖向奢侈、散漫、懒惰，并逐渐引向比这些更糟糕的恶果。因为在社会中，有许多情况能够扰乱心灵的平衡，阻碍其正直的进程；首先是与女性的社会交往。因为主教和照顾整个羊群的人不可能只关心男性部分，而忽略女性部分，而女性因易犯罪更需要特别照顾。但是被任命管理主教职的人必须关心她们的道德健康，如果不是更多，至少不少于其他人。因为有必要在她们生病时探访她们，在她们忧伤时安慰她们，在她们懒惰时责备她们，在她们痛苦时帮助她们；在这种情况下，恶者会找到许多接近的机会，如果一个人不用非常严格的守卫来巩固自己。因为不仅是淫荡的女人的眼睛，就是端庄的女人的眼睛也会刺透和扰乱心思。奉承使它疲弱，恩惠使它受奴役，而炽热的爱——可以说是一切善的源泉——对那些不善用它的人来说成了无数恶的原因。不断的忧虑也已经磨钝了理解的锐气，使轻快变得比铅还重，而愤怒像烟一样冲进来，占据了整个内在的人。\par

9. 何必谈论由忧伤带来的伤害、侮辱、辱骂、来自上者和下者、来自智者和愚者的谴责？因为那些缺乏正确判断的人特别喜欢谴责，并且从不轻易接受任何借口。但真正卓越的主教不应当轻视这些，而应当以极大的容忍和温良，向所有人澄清他们对他的指控，宽恕他们无理的挑剔，而不是对此愤慨和生气。因为如果圣保禄害怕在门徒中引起偷窃的嫌疑，因此让人管理钱财，他说：\Quote{免得有人在这丰裕的供给中责备我们}，\BibleRef{格后 8:20} 我们岂不应该尽一切努力消除恶意的猜疑，即使这些猜疑是假的、最不合理的、并且与我们想法相去甚远？因为我们没有一个人像圣保禄远离偷窃那样完全远离任何罪；尽管如此，他虽然远离这种恶行，却并没有因此轻视世人的猜疑，尽管这猜疑非常荒谬甚至疯狂。因为对那位有福和可钦佩的人产生任何这样的猜疑都是疯狂的。但他仍然远远地清除了这种猜疑的原因，尽管它不合理，并且没有任何头脑清醒的人会持有，他既没有藐视众人的愚蠢，也没有说：\Quote{谁曾想到在我们行了这些奇迹、过着如此克制的生活、并且你们都尊敬和钦佩我们之后，还会对我们产生这样的怀疑？} 恰恰相反，他预见到了这个卑鄙的猜疑，并将其连根拔起，或者更确切地说，根本没有让它生长。为什么？他说：\Quote{因为我们不仅在主面前，也在众人面前，预备诚实的事。} 因此，我们必须用同样的、甚至更大的热忱来根除和阻止不良的传言，并远远预见它们从哪里来，事先消除产生它们的原因，而不是等到它们被确立并成为人人谈论的话题。因为那时将来很难破坏它们，甚至可能不可能，而且不是没有危害，因为这是在许多人受到伤害之后才做的。但我还要继续追逐不可及的事情多久呢？因为列举这一方向的所有困难，无异于测量海洋。即使一个人使自己摆脱一切情感（这是不可能的），为了纠正他人的过失，他被迫承受无数的考验；而当他自己软弱的因素被加进来时，请看，劳苦和忧虑的深渊，以及所有希望克服自己和周围人罪恶的人必须忍受的一切。\par

10. 巴西略：而现在，你难道没有劳苦吗？你独自生活时就没有忧虑吗？\par

\Person{金口圣若望}：我的确现在也是如此。因为一个人既是人，又过着我们这种辛劳的尘世生活，怎能没有忧虑和争战呢？但是，一个人沉入无边的海洋与渡过一条河，两者并不完全相同，这些忧虑与那些忧虑之间的差别如此之大。因为现在，如果我能有益于他人，我自己也愿意这样做，这将是我祈求之事。但如果帮助别人不可能，那么只要能从波涛中拯救自己，我也就知足了。\par

\Person{巴西略}：那么你以为这是了不起的事吗？你以为当你不能使任何人受益时，你就能得救吗？\par

\Person{金口圣若望}：你说得好，而且高尚，因为我自己无法相信，一个人若不努力拯救他的邻人，还能得救；那福音中可怜的人也不曾因未减少他的塔冷通而受益，却因没有增加并加倍带给主人而丧亡。\BibleRef{玛 25:24} 然而，我觉得当我被算账时，我的惩罚会更轻，因为我未曾拯救别人，胜过我因在如此大的荣耀之后变得更加恶劣，而既毁了自己又毁了他人。因为我深信，我的惩罚将与我的罪孽相称；但在接受这职务后，我害怕惩罚将不只是双倍或三倍，而是多倍，因为我本会使许多人跌倒，并在接受了更大的荣耀后，冒犯了那荣耀我的天主。\par

为此天主更严厉地谴责以色列子民，并表明他们应受更重的惩罚，因为他们在从祂那里获得如此多荣耀之后犯罪，在一处说：\Quote{但我只从地上万族中认识了你们，因此我要惩罚你们的一切罪过。}\BibleRef{亚 3:2} 又说：\Quote{我从你们的子孙中兴起先知，从你们的青年中兴起\OriginalTerm{纳齐尔人}{Nazarites}。}\BibleRef{亚 2:11} 在先知时代以前，祂想要表明，当罪过发生在司铎身上时，远比发生在平信徒身上受到更严厉的惩罚，祂命令为司铎献上的祭品与为全民献上的一样多，这等于祂亲自证明，司铎职的创伤需要更多的医治——即与全体百姓的创伤一样大；它们本不需要更大，除非它们更严重；它们在本性上并非更严重，但因司铎的尊严而加重，因为他竟敢犯这些罪。我又何必提那些从事这种职务的\Emph{人}呢。因为司铎的女儿们，\BibleRef{肋 21:9} 她们在司铎职务上无分，却因她们父亲的尊严，为同样的罪受到比他人更严厉得多的惩罚；她们的过犯与平信徒女儿的过犯相同，即两者都是奸淫；但前者的刑罚远为严厉。你没有看见天主用多么充足的证据向你表明，祂要求统治者比被统治者受更大的惩罚吗？因为无疑地，一个人因另一个人的缘故而比别人更严厉地惩罚那人的女儿，他也不会向那导致她受额外惩罚的人索取与别人相同的惩罚，而是更重的惩罚；这非常合理，因为祸害不仅涉及他自己，还摧毁了软弱的弟兄和仰望他之人的灵魂。\Person{厄则克耳}为了表明这一点，在写作中区分了公绵羊和绵羊的审判。\BibleRef{则 34:17}\par

12. 那么我们这样看来，你们是否认为我们怀着合理的畏惧？因为除了已经说过的话，尽管我这边需要许多劳苦，以免完全被灵魂的情欲所淹没，但我还是忍受劳苦，并不逃避争战。因为即使现在我也被虚荣俘虏，但我常常恢复自我，一眼看出自己已被俘获，有时我斥责已被奴役的灵魂；甚至连无节制的欲望现在也临到我身上，但它们只燃起微弱的火焰，因为我肉眼无法找到任何燃料来喂养那火。至于说任何人的坏话，或听人谈论他人的恶事，我完全远离，因为我没有与人交谈的对象；因为这些墙壁绝不会开口说话。然而，避免愤怒并非完全可能，尽管没有人激起它。因为常常当我回忆起那些蛮横之人以及他们所行之事，这种回忆便使我心中愤慨。但不是永久的，因为我迅速扑灭它的燃起，并劝说它安静下来，说：\Quote{忽略自己的过失而忙于邻人的过失，是非常不明智且极为可鄙的。}但如果我进入人群，陷入无数的刺激中，我就不能受益于这种警告，也无法体验到这种自我谴责的反思。相反，正如那些被急流冲下悬崖或别样坠落的人，能够预见他们终将面临的毁灭，却无法想到任何救助之法，同样，当我陷入自己情欲的巨大喧嚣中，我将能够一眼看出我的惩罚日益增加。但要像我如今这样能自制，并谴责这种在四周肆虐的疾病，对我来说不会像以前那么容易。因为我的灵魂软弱而渺小，不仅容易受这些情欲的支配，也容易受比它们都更苦毒的嫉妒的支配。它也不知道如何适中地承受侮辱或荣誉。但后者过分地抬高它，而前者则压抑它。所以，正如凶猛的野兽，当它们状态良好、精力充沛时，会战胜与它们搏斗的人，尤其如果那些人软弱而不熟练；但如果有人通过饥饿削弱它们，就会使其狂怒休眠，并熄灭其大部分力量；这样，一个并非十分勇猛的人也可以起来与它们作战和搏斗；灵魂的情欲也是如此。使它们软弱的人，便使它们服从于正确的理性；而精心喂养它们的人，使他的争战更加艰难，使它们变得如此可畏，以至于他一生都处于奴役和恐惧之中。\par

那么，这些野兽的食物是什么呢？虚荣的食物确是荣誉和称赞；骄傲的食物是权势和权力的充裕；嫉妒的食物是邻人的名声；贪婪的食物是慷慨者的馈赠；不节制的食物是奢侈和与妇人不断的交往；其他情欲也有其适当的营养吗？如果我要进入世界，所有这些都会猛烈攻击我，将我的灵魂撕成碎片，它们将变得更加可畏，使我的争战更加艰难。而我安于此地时，它们将被制服；虽然只是以极大的努力，但同时，靠着天主的恩宠，它们\Emph{会}被制服，那时除了它们的吠叫，再没有更糟的东西了。由于这些原因，我留在这斗室里，不与人来往，自足而不合群，忍受听到无数这类抱怨，尽管我很乐意消除它们，并且因为不能做到而烦恼悲伤；因为对我来说，变得合群同时又保持目前的安全并不容易。因此，我也恳求你们，与其责备一个如此受困于巨大困难的人，不如怜悯他。\par

但我们还不能说服你们。因此，现在是我向你们说出我唯一没有说出口的事情的时候了。也许在许多看来似乎难以置信，但即便如此，我也不羞于将其公之于世，因为虽说所说的是邪恶良心和许多罪的证据，但既然将要审判我们的天主知道一切细节，那么从人的无知中我们能获得什么益处呢？那么，这尚未说出口的是什么呢？从你们向我暗示主教职的那一天起，我的整个系统常常处于完全崩溃的危险中，如此大的恐惧，如此大的沮丧攫住了我的灵魂；因为一思想基督的新娘的荣耀、圣洁、属灵的美貌与智慧、及秀美，然后数算我自己的过失，我便不断哀哭她和自己，在不断的痛苦和困惑中，我一直在说——\Quote{那么是谁提出了这样的建议？为什么天主的教会犯了如此大的错误？为什么她如此激怒她的主人，以致被交给我，最不配的所有人，并经受如此大的耻辱？}我常常独自思量这些事，无法忍受如此可怕之事的念头，我就像被雷击的人一样，哑口无言，既不能看也不能听。当这种极度无助的状态离开我时——有时它会过去——泪水和沮丧接踵而来，在泪潮之后，恐惧又取而代之，扰乱、迷惑、搅动我的心神。在这样的风暴中，我度过了过去的时光；但你们却不知道，以为我过着完全平静的生活，但现在我将努力向你们揭示我灵魂的风暴，因为或许你们会从此原谅我，放弃你们的指控。那么，我该如何向你们揭示呢？因为如果你们要清楚地看到这个，除了袒露自己的心之外别无他法；但既然这是不可能的，我将尽力通过一个模糊的比喻向你们展示我沮丧的阴暗，并请从这个图像推断我的状况。\par

让我们设想，普天之下万王之王的女儿是某人的未婚妻，而这少女有着无与伦比的美貌，超越凡人的本性，在这方面她远远超过了全体妇女；此外，她拥有如此伟大的灵魂美德，以至于远远超越了所有过去和将来的人；她的举止恩宠超越了所有艺术的标准，而她的容貌之美使得她本人的可爱黯然失色；她的未婚夫不仅为了这些而爱恋这少女，而且除此之外还对她怀有深情，他的热情使曾经最炽热的情人都相形见绌。那么，让我们设想，当他燃烧着爱情时，却从某处听说，一个卑贱、低劣的人，出身低微、身体残疾，实际上是一个彻头彻尾的恶棍，将要迎娶这位奇妙而深爱的少女。我们是否已经向你呈现了我们悲伤的一小部分？我的比喻到此为止是否足够？就我的沮丧而言，我认为已经够了；因为我引入这个比喻的唯一目的就是此，但为了向你展示我的恐惧和惊骇的程度，让我继续另一个描述。\par

让我们设想一支由步兵、骑兵和海军组成的军队，让许多三层桨战船覆盖海面，步兵和骑兵的方阵覆盖大部分平原和山脊，他们的盔甲在阳光下反射光芒，头盔和盾牌的闪光被射出的光线反射；长矛的碰撞声和马匹的嘶鸣声直冲云霄，无论海洋还是陆地都看不见，只有四面八方的铜和铁。让敌军列阵相对，他们是凶猛野蛮的人，交战时刻即将来临。然后，让某人突然抓住一个年轻的小伙子，一个在乡间长大、只知道牧羊笛和牧杖的人；给他穿上铜甲，带他绕整个营地，让他看到中队和他们的军官、弓箭手、投石手、队长、将军、步兵和骑兵、长矛手、三层桨战船及其指挥官、船上密集的士兵以及船上准备好的大量战争器械。此外，让他看到敌人的整个阵列，他们可怖的面容，以及他们武器多样且不寻常的数量；还有山中的沟壑和悬崖，深邃而险峻。再让他看到敌军方面，马匹如被魔法般飞行，步兵在空中行进，以及各种威力与形式的巫术；让他考虑战争的灾难，长矛的云层，箭矢的冰雹，那巨大的迷雾和黑暗，那由大量武器引起的、用它们的云层遮蔽阳光的最黑暗的夜晚，灰尘和黑暗同样使眼睛模糊。血流的洪流，倒下者的呻吟，幸存者的呼喊，被杀者的堆叠，被血浸透的车轮，因尸体众多而连人带马摔下的骑手，地面一片混乱，血、弓、箭、马蹄和人头混杂在一起，人的手臂、战车轮子和头盔，被刺穿的胸膛，粘在剑上的脑浆，折断的矛尖上还插着一只眼睛。然后让他估算海军的苦难，三层桨战船在波涛中燃烧，连同武装的船员一起沉没，海的咆哮，水手的骚动，士兵的呼喊，混着血的波浪泡沫，溅入所有船只；甲板上的尸体，有些下沉，有些漂浮，有些被抛上海滩，被波浪淹没，阻碍船只航行。当他仔细了解了战争的全部悲剧之后，再加上俘虏和奴役的恐怖，这比任何死亡都更可怕；告诉他这一切之后，命令他立即上马，指挥那整个军队。\par

你真的认为这个少年能够承受这仅仅的描述，而不会在最初一瞥时就失去勇气吗？\par

13. 不要以为我夸张了此事，也不要因为我们是困在这身体里如同在监牢中，无法看见那看不见的世界，就以为我所说的是夸大其词。因为倘若你能用你这肉眼看见魔鬼那最阴暗的阵势及其疯狂的进攻，你必会看到一场比这更大更可怕的战斗。那里没有铜铁，没有马匹、战车或车轮，也没有火与箭。这些是可见之物。但有比这些更可怕得多的兵器。对抗这些敌人并不需要胸甲或盾牌、刀剑或长矛，然而仅仅是这受诅咒的阵列的景象就足以使灵魂瘫痪，除非这灵魂极其高贵，并且在天主的眷顾中享有高度的保护作为其勇气的庇护。如果可能，脱下这身体，或仍留着它，用肉眼清晰无畏地看见他的整个阵势以及他对我们的战争，你将不会看见血流成河或尸体，而是看见许多堕落的灵魂和如此悲惨的创伤，以至于我刚才详细描述的那种战争，你会认为不过是儿戏和消遣，而非战争。每天有如此多人被击中，而这两种情况下的创伤并不导致同样的死亡；灵魂与身体的差异有多大，那种死亡与这种死亡的差异就有多大。因为当灵魂受伤并倒下时，它并不像一具无生命的身体那样躺卧，而是从此受折磨，被邪恶的良心啃噬；并且在此世离世之后，在审判之时，它被交付于永罚；如果有人对魔鬼所给的创伤毫不悲伤，他的危险便因麻木而更大。因为谁不为第一次创伤感到痛苦，就会轻易接受第二次，接着第三次。因为那污秽的灵不会停止攻击，直到最后一息，每当他发现一个灵魂对最初的创伤漫不经心、漠不关心时；如果你询问其攻击的方法，你会发现这方法更加猛烈多样。因为没有人曾知道如此多的诡计和欺骗的形式，如同那污秽的灵。正因如此，它获得了大部分力量，也没有人能像那恶者对待人类那样，对最坏的敌人怀有如此不可和解的仇恨。如果有人询问它战斗的猛烈程度，同样，将人与它相比也是可笑的。但若有人选出最凶猛、最野蛮的野兽，想要将其狂暴与它相比，他会发现那些野兽相比之下是温顺安静的，因为它在攻击我们的灵魂时散发出如此的怒气；前者的战争时期短暂，并且在这短暂的空间里有间歇；因为黑夜的来临、杀戮的疲劳、用餐时间以及许多其他事情，给士兵以喘息，使他能卸下盔甲，稍作呼吸，用食物和饮料恢复精力，并以许多其他方式恢复他原有的力量。但在恶者的情况下，永远不可能卸下盔甲，甚至连睡眠也不可能，对于那些想要始终安然无恙的人。因为二者必居其一：要么倒下赤手空拳而灭亡，要么全副武装且时刻警觉。因为它总是带着自己的阵势站立，窥伺我们的懈怠，为我们的毁灭比我们为自己的救恩更加热心地劳作。\par

它不被我们看见，却突然袭击我们——这对那些不常警觉的人来说是无数祸患的根源——这也证明这种战争比另一种更难。你难道希望我们在这种情况下指挥基督的士兵吗？是的，这等于把他们指挥给魔鬼服务，因为每当那本应列阵并指挥他人的人是最无经验、最软弱的人时，他通过这种无经验背叛那些托付给他的人，就等于是为魔鬼的利益而非基督的利益指挥他们。\par

但你为什么叹息？为什么哭泣？因为我的安逸如今并不需要哀哭，而是需要喜乐和欢欣。\par

巴希尔：但我的情况并非如此，是的，这需要无数的哀哭。因为我几乎还不能明白你把我带到了何等邪恶的境地。因为我来到你这里，想要知道在那些指责你的人面前，我该为你作出怎样的辩护；但你却把另一个案件放在我原来的案件之后，打发我回去。因为我不再关心我该在他们面前为你给出什么辩护，而是我该为自己和自己的过失向天主作出什么辩护。但我恳求你，并求你，如果我的福利在你眼中还有价值，如果在基督内有什么劝勉，如果有什么爱的安慰，如果有什么慈悲和怜悯（\BibleRef{斐 2:1}），因为你知道是你首先把我带入这危险之中，求你伸出援手，说出并做出能够恢复我的事，不要忍心在最短的时刻离开我，而要如今比以前更让我与你共度此生。\par

金口：我微笑着说道，在如此重大的职务重担之下，我怎能帮助你，怎能有益于你呢？但既然这对你是愉快的，请鼓起勇气，亲爱的灵魂，因为只要你在自己的挂虑中有空暇，我就会来安慰你，在我的能力范围内，没有一样会缺少。\par

听到这话，他更加哭泣，站起身来。但我拥抱他，亲吻他的头，领他出去，劝他勇敢地承受他的命运。因为我相信，我说，藉着那召叫了你并把你安置在他自己羊群之上的基督，你会从这职务中获得如此的确信，以至于在末日，如果我处于危险之中，你也会将我接到你永恒的帐幕里。\par

\SourceRecord{Translated by W.R.W. Stephens.；Nicene and Post-Nicene Fathers, First Series；Vol. 9.；Edited by Philip Schaff.；Buffalo, NY: Christian Literature Publishing Co.,；1889.；<http://www.newadvent.org/fathers/19226.htm>.}
